\section{Causal Relations and Confounding---Leftovers}

\subsection{Direct confounding}

\begin{Def}\label{def:graph_directly_confounded}
  Let $M = \lt \Sin, \Send, \Srnd, \CX, P, f \rt$ be a simple iSCM with causal graph $G(M)$.
  Let $a,b \in V$ be distinct endogenous variables.
  If $a \huh b$ is present in $G(M)$ then we say that \emph{$a$ and $b$ are directly confounded (w.r.t.\ $V \cup J$) according to $M$}.
\end{Def}
\begin{Rem}\label{rem:graph_directly_confounded}
  In terms of the graph $G^+(M)$ of $M$, this holds if and only if there exists $w \in W$ such that $a \hut w \tuh b$ in $G^+(M)$.
\end{Rem}
\begin{Def}\label{def:scm_directly_confounded}
Let $M = \lt \Sin, \Send, \Srnd, \CX, P, f \rt$ be a simple iSCM.
Let $a, b \in V$ with $a \ne b$.
If for all values $x_J \in \CX_J, x_V \in \CX_V$, $f_a(x_J,X_W,x_V) \Indep f_b(x_J,X_W,x_V)$ then we say that \emph{$a$ and $b$ are ``directly'' unconfounded according to $M$}.
  \Joris{This might work in the acyclic case, but in the cyclic case we really have to solve $M^{[a,b]}$.}
\end{Def}
\begin{Prp}\label{prp:directly_confounded}
  Let $M = \lt \Sin, \Send, \Srnd, \CX, P, f \rt$ be a simple iSCM with causal graph $G(M)$.
  For $a,b \in V$, $a$ and $b$ are directly confounded (w.r.t.\ $V \cup J$) according to $M$ if and only if $a$ and $b$ are confounded according to $M^{[\{a,b\}]}$.
  \Joris{??? update!}
\end{Prp}

\begin{Prp}\label{prp:scm_unconfoundedness}
Let $M = \lt \Sin, \Send, \Srnd, \CX, P, f \rt$ be a simple iSCM.
Let $a \ne b \in V$.
If $b \notin \Anc^{G(M)}(a)$ and $a$ and $b$ are not confounded according to $M$, then
  $$P_{M}(X_b \given \doit(X_a), \doit(X_{V\setminus\{a,b\}}), \doit(X_J)) = P_{M}(X_b \given X_a, \doit(X_{V\setminus\{a,b\}}), \doit(X_J)) \qquad \mu_a\text{-a.s.}$$
for any reference measure $\mu_a$ on $\CX_a$ with $\mu_a \ll P_M(X_a \given \doit(X_J)) \ll \mu_a$.
\end{Prp}
%The contra-positive can be used as a sufficient condition for detecting that two nodes are confounded according to $M$.

For the simpler representation, we can intervene on all variables in $V\sm\{a,b\}$ and then only a few possible edges remain: $a \tuh b$, $b \tuh a$, $a \huh b$. 
If $a \huh b$ isn't there, we get the separation assumption and can draw the conclusion of no ``direct'' confounding bias (if we intervene to fix all other variables).

This leads to a criterion to detect a bidirected edge!


%Note that exogenous input nodes in $J$ are never called confounders, by definition.\footnote{The reason is that they do not lead to ``spurious dependences'' as long as one does not mix distributions for different values of $X_J$. 
%\Joris{How to deal with constant exogenous random nodes, or constant endogenous nodes without parents? We might want to not consider those as confounders. Also, the Markov property could be extended: by substituting the constant values in their children nodes (perhaps recursively) we can get rid of these nodes, apply the Markov property to the result, and we can freely add the removed nodes in any of the spots of any separation statement.}}
\subsection{Faithfulness}

\Joris{TODO: add an example of a four-cycle that is linear and hence $d$-faithful but not $\sigma$-faithful.}
  \begin{Eg}\label{eg:simple_scm_d_faithful_not_sigma_faithful}
  \begin{figure}\centering
  \begin{tikzpicture}
      \node[ndout] (X1) at (0,0) {$X_1$};
      \node[ndout] (X2) at (1.5,0) {$X_2$};
      \node[ndout] (X3) at (1.5,1.5) {$X_3$};
      \node[ndout] (X4) at (0,1.5) {$X_4$};
      \node[ndlat] (W1) at (-1,-1) {$W_1$};
      \node[ndlat] (W2) at (2.5,-1) {$W_2$};
      \node[ndlat] (W3) at (2.5,2.5) {$W_3$};
      \node[ndlat] (W4) at (-1,2.5) {$W_4$};
      \draw[arout] (X1) edge (X2);
      \draw[arout] (X2) edge (X3);
      \draw[arout] (X3) edge (X4);
      \draw[arout] (X4) edge (X1);
      \draw[arout] (W1) edge (X1);
      \draw[arout] (W2) edge (X2);
      \draw[arout] (W3) edge (X3);
      \draw[arout] (W4) edge (X4);
  \end{tikzpicture}
  \caption{Graph $G^+(M)$ of SCM $M$ of Example~\ref{eg:simple_scm_d_faithful_not_sigma_faithful}. 
  $X_1 \dPerp_{G^+(M)} X_3 \given X_2, X_4$, but $X_1 \sPerp_{G^+(M)} X_3 \given X_2, X_4$.
  $M$ is $d$-faithful, but not $\sigma$-faithful.}
  \end{figure}
Consider an SCM $M$ with four endogenous variables $X_1,X_2,X_3,X_4$, four exogenous random variables $W_1,W_2,W_3,W_4$, with structural equations
  \begin{align*}
    X_1&=X_4+W_1 \\
    X_2&=X_1 \cdot W_2 \\
    X_3&=X_2+W_3 \\
    X_4&=X_3\cdot W_4
  \end{align*}

Consider an SCM with structural equations
  \begin{align*}
    X_1 &= W_1 \\
    X_2 &= W_2 \\
    X_3 &= X_1 + X_4 + W_3 \\
    X_4 &= X_2 + X_3 + W_4 
  \end{align*}
where the $X$'s are considered real-valued endogenous variables and the $W$'s exogenous variables with domains $(-1,1) \subset \RN$.
We can solve the system of structural equations for $X$ in terms of $W$:
  \begin{align*}
    X_1 &= W_1 \\
    X_2 &= W_2 \\
    X_3 &= \frac{W_3 + W_1 W_4}{1 - W_1 W_2} \\
    X_4 &= \frac{W_2 W_3 + W_4}{1 - W_2 W_1}
  \end{align*}
\end{Eg}

The following:
$$X_A \Indep_{P_{M}(V,W|J)} X_B \given X_C,$$
implies that for every probability distribution $Q(J) \in \Pcal(\CXJ)$ we have the conditional independence:
\[ X_A \Indep_{P_{M}(V,W|J)\otimes Q(J)} X_B,X_J \given X_C\]

We can now make the following faithfulness assumption:
assume that
\[ X_A \Indep_{P_{M}(V,W|J)\otimes Q(J)} X_B, X_J \given X_C\]
for all $Q(J) \in \Pcal(\CXJ)$ 
implies
\[ A \isPerp_{G(M)} B \given C\]
By the Markov property, then, we conclude
  $$X_A \Indep_{P_{M}(X_V,X_W \given \doit(\XJ))} X_B \given X_C$$

Thus under that faithfulness assumption, we have
\[ X_A \Indep_{P_{M}(V,W|J)\otimes Q(J)} X_B,X_J \given X_C\]
for all $Q(J) \in \Pcal(\CXJ)$ if and only if
\[ A \isPerp_{G(M)} B \given C\]
if and only if
\[ A \isPerp_{\hat{G}(M)} B \cup J \given C\]
where $\hat{G}(M)$ is like $G(M)$ but with bidirected edges between all input nodes,
and then input nodes replaced by normal nodes. This means that we have reformulated 
everything in terms of older concepts. However, to arrive at this conclusion, our long
excursion was necessary. And under this particular faithfulness assumption, all statements
in Theorem~\ref{thm-ci-equivalences} are equivalent, and also equivalent to the 
corresponding $i\sigma$-separation.

With the newer concepts, we can alternatively assume
  $$X_A \Indep_{P_{M}(X_V,X_W \given \doit(\XJ))} X_B \given X_C$$
implies
\[ A \isPerp_{G(M)} B \given C\]

But I think we need something else. It would be nice if we could assume that if
\[ X_A \Indep_{P_{M}(V,W|J)\otimes Q(J)} X_B, X_J \given X_C\]
for \emph{some} $Q(J) \in \Pcal(\CXJ)$, then
\[ A \isPerp_{G(M)} B \given C\]

\subsection{Causal Relations---Separate Cases}

We are now ready to formalize Definition~\ref{def-causality-real} (finally!), that is, the notion of ``$a$ causes $b$''.
%In this chapter, we will frequently be making use of the following short-hand notation.
%\begin{Def}
%  Let $G$ be a CDMG with input nodes $J$ and output nodes $V$.
%  For $O \subseteq V$, we write $G_{O|J} := G^{\sm (V \sm O)}$ for the graph obtained by marginalizing all output nodes except for those in $O$.
%\end{Def}

\subsubsection{Causal relations (graphical notion)}

\begin{Def}\label{def:graph_cause_of}
  Let $G$ be a CDMG with input nodes $J$ and output nodes $V$.
  Let $a \in J \cup V$ and $b \in V$.
  When interpreting the graph causally, we say that \emph{$a$ is a cause of $b$ according to $G$} if $a \in \Anc^G(b)$.
\end{Def}
\begin{Rem}\label{rem:graph_cause_of}
  Marginalization of a graph preserves its causal relations (see also Remark~\ref{rem:marg_preserves_ancestors_bifurcations_acyclicity}).
  In particular, $a$ is a cause of $b$ according to $G$ if and only if $a \tuh b$ is present in in $G^{\sm(V\sm\{a,b\})}$.
  Furthermore, $a$ is a cause of $b$ according to $G$ if and only if $a$ is a cause of $b$ according to $G_{\doit(a)}$.
\end{Rem}
\begin{Prp}\label{prp:graph_cause_of}
  Let $G$ be a CDMG with input nodes $J$ and output nodes $V$.
  Let $b \in V$.
  For $a \in J \cup V$:
    $$a \not\in \Anc^G(b) \iff b \isPerp_{G_{\doit(a)}} a \mid J \setminus \{a\}.$$
  Hence, for $a \in V$:
    $$a \not\in \Anc^G(b) \iff b \isPerp_{G_{\doit(I_a)}} I_a \mid J.$$
\end{Prp}
\begin{proof}
  Suppose that there exists a walk in $G_{\doit(a)}$ between $b$ and $a \cup J$ that is $\sigma$-open given $J \setminus \{a\}$.
  Then there exists such a walk without colliders (Proposition~\ref{prp:sigma_opens} and the observation that no node in $J$ can be a collider on a walk).
  Such a walk cannot contain a node from $J \setminus \{a\}$, since that would either be an end node ($\sigma$-blocking the walk) or a non-collider node pointing only to nodes in another strongly connected component ($\sigma$-blocking the walk).
  Therefore it must be of the form $a \tuh \cdots b$.
  Since it cannot contain a collider, it must be a directed walk, and hence $a \in \Anc^G(b)$.
  Vice versa, any directed walk from $a$ to $b$ in $G_{\doit(a)}$ is $\sigma$-open given $J \setminus \{a\}$, as it cannot contain a node in $J \sm \{a\}$.

  Suppose $a \in V$.
  Then $a \in \Anc^{G}(b)$ holds if and only if $I_a \in \Anc^{G_{\doit(I_a)}}(b)$ holds, which holds if and only if
    $$b \isPerp_{G_{\doit(I_a)}} I_a \mid J$$
  according to the first statement applied to $G_{\doit(I_a)}$, $I_a$ and $b$.
%  Then, if there exists a walk in $G_{\doit(a)}$ between $b$ and $a \cup J$ that is $\sigma$-open given $J \setminus \{a\}$, it must be of the form $a \tuh \cdots b$ (as we just saw).
%  It must also be present in $G_{\doit(I_a)}$, and we can extend it to $I_a \tuh a \tuh \cdots b$, which gives a walk in $G_{\doit(I_a)}$ that is $\sigma$-open given $J$.
%  Vice versa, if there exists a walk in $G_{\doit(I_a)}$ between $b$ and $I_a \cup J$ that is $\sigma$-open given $J$, then it must be of the form $I_a \tuh a \tuh \dots b$, as it cannot contain a node from $J$, and cannot contain a collider.
%  The subwalk $a \tuh \dots b$ is then $\sigma$-open given $J$ in $G$.
%  Hence there must exist a path $a \tuh \dots b$ in $G$ that is $\sigma$-open given $J$.
%  This must also exist in $G_{\doit(a)}$ and must be $\sigma$-open given $J$.
%  Hence there exists a path in $G_{\doit(a)}$ between $b$ and $a$ that is $\sigma$-open given $J \sm \{a\}$.
\end{proof}
\Joris{
\begin{Rem}
  If we would identify $I_a$ with $a$ for $a \in J$, then we could also combine the two cases as:
    $$b \isPerp_{G_{\doit(I_a)}} I_a \mid J \sm \{I_a\}.$$
  for $a \in J \cup V$.
  (Note that for $a \in J \cup V$, $a \in \Anc^G(b) \iff I_a \in \Anc^{G_{\doit(I_a)}}(b)$).
\end{Rem}
}


\Joris{
  What if we identify $I_a$ with $a$ for $a \in J$?
  Suppose we consider, for $c \in J \cup I_V$:
    $$c \not\in \Anc^G(b) \iff b \isPerp_{G_{\doit(c)}} c \mid J \sm \{c\}.$$
  Specializing to $c = a \in J$:
    $$a \not\in \Anc^G(b) \iff b \isPerp_{G} a \mid J \sm \{a\}.$$
  Specializing to $c = I_a \in I_V$:
    $$I_a \not\in \Anc^G(b) \iff b \isPerp_{G_{\doit(I_a)}} I_a \mid J.$$
  So it works!

UNIFICATION:
\begin{Prp}\label{prp:graph_cause_of}
  Let $G$ be a CDMG with input nodes $J$ and output nodes $V$.
  Let $b \in V$.
  For $a \in J \cup V$:
    $$a \not\in \Anc^G(b) \iff b \isPerp_{G_{\doit(I_a)}} I_a \mid J \sm \{I_a\}.$$
\end{Prp}
\begin{proof}
  Suppose that there exists a walk in $G_{\doit(I_a)}$ between $b$ and $I_a \cup J$ that is $\sigma$-open given $J \setminus \{I_a\}$.
  Then there exists such a walk without colliders (Proposition~\ref{prp:sigma_opens} and the observation that no node in $J$ can be a collider on a walk).
  Such a walk cannot contain a node from $J \setminus \{I_a\}$, since that would either be an end node ($\sigma$-blocking the walk) or a non-collider node pointing only to nodes in another strongly connected component ($\sigma$-blocking the walk).
  Therefore it must be of the form $I_a \tuh a \tuh \cdots b$.
  Since it cannot contain a collider, it must be a directed walk, and hence $I_a \in \Anc^{G_{\doit(I_a)}}(b)$.
  Vice versa, any directed walk from $I_a$ to $b$ in $G_{\doit(I_a)}$ is $\sigma$-open given $J \setminus \{I_a\}$, as it cannot contain a node in $J \sm \{I_a\}$.
  Finally, note that for $a \in J \cup V$, $a \in \Anc^G(b) \iff I_a \in \Anc^{G_{\doit(I_a)}}(b)$.
\end{proof}
}

\subsubsection{Causal relations (interventional notion)}
The most common notion of causation is the interventional one.
\begin{Def}\label{def:scm_rung2_cause_of}
  Let $M = \lt \Sin, \Send, \Srnd, \CX, P, f \rt$ be a simple iSCM.
  Let $a \in J \cup V$ and $b \in V$.
  We say that \emph{$a$ is a rung-2 cause of $b$ according to $M$} if, in case $a \in V$:
    $$P_{M}(X_b \given \doit(X_J), \doit(X_a)) \ne P_{M}(X_b \given \doit(X_J))$$
  and, in case $a \in J$:
    $$P_{M}(X_b \given \doit(X_J)) \ne P_{M}(X_b \given \doit(X_{J\setminus \{a\}}),\cancel{\doit(X_a)}).$$
\end{Def}
\begin{Rem}\label{rem:scm_rung2_cause_of}
  Note: both cases can also be combined as a single statement:
    $$P_{M}(X_b \given \doit(X_{J \cup \{a\}})) \ne P_{M}(X_b \given \doit(X_{J\sm\{a\}})).$$
  Also note that ``$a$ is not a rung-2 cause of $b$ according to $M$'' is equivalent to:
  \begin{equation*}
    \begin{cases}
      X_b \Indep_{P_{M}} X_a \given X_{J \sm \{a\}} & a \in J, \\
      X_b \Indep_{P_{M_{\doit(I_a)}}} X_{I_a} \given X_J & a \in V.
    \end{cases}
  \end{equation*}
  \Joris{Hence, if we would identify $I_a$ with $a$ for $a \in J$, then we could also combine the two cases as:
    $$X_b \Indep_{P_{M_{\doit(I_a)}}} X_{I_a} \given X_{J \sm \{I_a\}}.$$
    for $a \in J \cup V$.
    }
\end{Rem}
\Joris{
  What if we identify $I_a$ with $a$ for $a \in J$?
  Suppose we consider, for $c \in J \cup I_V$:
    $$X_b \Indep_{P_{M_{\doit(c)}}} X_c \given X_{J \sm \{c\}}.$$
  Specializing to $c = a \in J$:
    $$X_b \Indep_{P_{M}} X_a \given X_{J \sm \{a\}}.$$
  Specializing to $c = I_a \in I_V$:
    $$X_b \Indep_{P_{M_{\doit(I_a)}}} X_{I_a} \given X_J.$$
  So it works!

  UNIFICATION:

\begin{Def}\label{def:scm_rung2_cause_of}
  Let $M = \lt \Sin, \Send, \Srnd, \CX, P, f \rt$ be a simple iSCM.
  Let $a \in J \cup V$ and $b \in V$.
  We say that \emph{$a$ is a rung-2 cause of $b$ according to $M$} if:
    $$X_b \Indep_{P_{M_{\doit(I_a)}}} X_{I_a} \given X_{J \sm \{I_a\}}.$$
\end{Def}
Note: if $M$ and $\tilde{M}$ are interventionally equivalent, then so are $M_{\doit(I_a)}$ and $\tilde{M}_{\doit(I_a)}$ for $a \in (V \cap \tilde{V}) \cup (J \cap \tilde{J})$.
  }
\begin{Rem}\label{rem:scm_rung2_cause_of_marginalization}
  Marginalization of iSCMs preserve the rung-2 causal relations between the remaining variables.
\end{Rem}
\begin{Rem}\label{rem:scm_rung2_cause_of_interventional_equivalence}
  Rung-2 causal relations are invariant under interventional equivalence of iSCMs.
  More precisely: $a$ is a rung-2 cause of $b$ according to $M$ implies that $a$ is a rung-2 cause of $b$ according to $\tilde{M}$ if $M$ is interventionally equivalent to $\tilde{M}$ w.r.t.\ $\{a,b\} \cap V$.
\end{Rem}
Definition~\ref{def:scm_rung2_cause_of} looks slightly involved as it differentiates between the cases $a \in V$ and $a \in J$.
The reason that we do this is because for $a \in V$, we can also make use of information in the observational distribution, as the following example illustrates.
\Joris{Common terminology is a bit sloppy: if we say `treatment has no effect on outcome', then we could mean `the treatment variable (with values A/B)' has no effect,
or `the treatment variable (with values A/B/no treatment)' has no effect. So would it be better to simplify the mathematical definition, and realize that if the
latter is meant, we have to first extend the model with an intervention variable and then we can apply the definition to that `treatment variable'?}
\begin{Eg}\label{ex:scm_cause_subtlety}
  Consider the SCM $M$ with endogenous variables $X_1,X_2,X_3 \in \{-1,+1\}$, exogenous random variable $E \in \{-1,1\}$, and structural equations:
  $$\begin{aligned}
      X_1 &= X_3 \\
      X_2 &= X_1 X_3 \\
      X_3 &= E
    \end{aligned}$$
  where $E \sim \C{U}(\{-1,1\})$.
  Its graph is shown in Figure~\ref{fig:scm_cause_subtlety}.
  Both interventional distributions $P_M(X_2\given\doit(X_1=-1)) = P_M(X_2\given\doit(X_1=1))$ coincide, but the marginal observational distribution $P_M(X_2)$ differs from both.\footnote{This example shows that the claim in \cite[Proposition 6.13]{PJS17} is incorrect.}
  \begin{figure}
  \begin{center}
    \begin{tikzpicture}
      \begin{scope}
        \node[ndout] (X1) at (0,0) {$X_1$};
        \node[ndout] (X2) at (2,0) {$X_2$};
        \node[ndout] (X3) at (1,1) {$X_3$};
        \node[ndlat] (E3) at (1,2.2) {$E$};
        \draw[arout] (X3) -- (X1);
        \draw[arout] (X1) -- (X2);
        \draw[arout] (X3) -- (X2);
        \draw[arout] (E3) -- (X3);
      \end{scope}
    \end{tikzpicture}
  \end{center}
  \caption{Graph $G^+(M)$ for the SCM $M$ of Example~\ref{ex:scm_cause_subtlety}.\label{fig:scm_cause_subtlety}}
\end{figure}
\end{Eg}


\subsubsection{Causal relations (counterfactual notion)}

A more `sensitive' notion of causation can be obtained by considering counterfactuals.
\begin{Def}\label{def:scm_rung3_cause_of}
  Let $M = \lt \Sin, \Send, \Srnd, \CX, P, f \rt$ be a simple iSCM.
  Let $a \in J \cup V$ and $b \in V$.
  We say that \emph{$a$ is a rung-3 cause of $b$ according to $M$} if, in case $a \in V$:
%    $$P_{M}(X_b \given \doit(X_J), \doit(X_a)) \ne P_{M}(X_b \given \doit(X_J))$$
    $$\exists x_J \in \CX_J, \exists x_a \in \CX_a: \quad P_{((M_{\doit(X_J=x_J)})^\twin)_{\doit(a)}}(X_b = X_{b'} \given \doit(X_a=x_a)) \ne 1;$$
  and, in case $a \in J$:
%    $$P_{M}(X_b \given \doit(X_J)) \ne P_{M}(X_b \given \doit(X_{J\setminus \{a\}}),\cancel{\doit(X_a)}).$$
    $$\exists x_{J\sm\{a\}} \in \CX_{J\sm\{a\}}, \exists x_a,x_a' \in \CX_a: \quad P_{(M_{\doit(X_{J\sm\{a\}}=x_{J\sm\{a\}})})^\twin}(X_b = X_{b'} \given \doit(X_a=x_a),\doit(X_{a'}=x_a')) \ne 1.$$
\end{Def}
\Joris{
  What if we identify $I_a$ with $a$ for $a \in J$?
  Can we then combine both cases into:
    \begin{equation*}\begin{split}
      & \exists x_{J\sm\{I_a\}} \in \CX_{J\sm\{I_a\}}, \exists x_{I_a},x_{I_{a'}} \in \CX_a: \\
      & P_{(M_{\doit(X_{J\sm\{I_a\}}=x_{J\sm\{I_a\}}),\doit(I_a)})^\twin}(X_b = X_{b'} \given \doit(X_{I_a}=x_{I_a}),\doit(X_{I_{a'}}=x_{I_{a'}})) \ne 1?
    \end{split}\end{equation*}
  Yes we can!
}
\begin{Rem}\label{rem:scm_rung3_cause_of_marginalization}
  Marginalization of iSCMs preserves the rung-3 causal relations of the remaining variables.
\end{Rem}
\begin{Rem}\label{rem:scm_rung3_cause_of_counterfactual_equivalence}
  Rung-3 causal relations are invariant under counterfactual equivalence of iSCMs.
  More precisely: $a$ is a rung-3 cause of $b$ according to $M$ implies that $a$ is a rung-3 cause of $b$ according to $\tilde{M}$ if $M$ is counterfactually equivalent to $\tilde{M}$ w.r.t.\ $\{a,b\} \cap V$.
\end{Rem}

The following example shows that the presence of a rung-3 causal relation does not necessarily imply the presence of a rung-2 causal relation.
Intuitively, the causal effect can go unnoticed if it averages out.
\begin{Eg}\label{eg:scm_cause_of_rung2_vs_rung3}
  Consider the iSCM $M$ with exogenous input variable $X_a \in \{-1,+1\}$, endogenous variable $X_b \in \{-1,+1\}$, exogenous random variable $X_w \in \{-1,1\}$, and structural equation:
  $$X_b = X_a X_w,$$
  where $X_w \sim \C{U}(\{-1,1\})$.
  Note that $a$ is a rung-3 cause of $b$ according to $M$, but $a$ is not a rung-2 cause of $b$ according to $M$.
  Indeed, due to the uniform distribution of $X_w$ the dependence of $b$ on $a$ averages out.
\end{Eg}

\subsubsection{Hierarchy of causal relations}
The following result shows the logical relations between various notions of causation.
We first consider the case $a \in J$, and use that to deal with the other case $a \in V$.
\begin{Prp}\label{prp:scm_exogenous_input_causes}
  Let $M = \lt J, V, W, \CX, P, f \rt$ be a simple iSCM with causal graph $G(M)$.
  Let $a \in J$, $b \in V$. 
  Pick an arbitrary $x_a^0 \in \CX_a$.
  Write ``$\foralmostall x_w \in \CX_W$'' as a shorthand of ``for $P$-almost all $x_w \in \CX_W$''.
  Denote the solution function of $M$ by $g$.
  Consider the properties:\footnote{Note that the ordering of the quantifiers matters! In particular, $\foralmostall x_W \forall x_V \implies \forall x_V \foralmostall x_W$, whereas the converse doesn't hold in general (but does hold if $\CX_V$ is discrete).}
  \begin{enumerate}[label=(\roman*)]
    \item[(i)]   $a \not\in \Anc^{G(M)}(b)$;
    \item[(ii)]  $\forall x_{J\sm\{a\}} \in \CX_{J\sm\{a\}}$, $\forall x_W \in \CX_W$, $\forall x_a \in \CX_a$: $g_b(x_a,x_{J\sm\{a\}},x_W) = g_b(x_a^0,x_{J\sm\{a\}},x_W)$;
%    \item[(iii)] $\foralmostall x_W \in \CX_W$, $\forall x_{J\sm\{a\}} \in \CX_{J\sm\{a\}}$, $\forall x_a \in \CX_a$: $g_b(x_a,x_{J\sm\{a\}},x_W) = g_b(x_a^0,x_{J\sm\{a\}},x_W)$;
    \item[(iii)]  $\forall x_{J\sm\{a\}} \in \CX_{J\sm\{a\}}$, $\foralmostall x_W \in \CX_W$, $\forall x_a \in \CX_a$: $g_b(x_a,x_{J\sm\{a\}},x_W) = g_b(x_a^0,x_{J\sm\{a\}},x_W)$;
    \item[(iv)]   $\forall x_{J\sm\{a\}} \in \CX_{J\sm\{a\}}$, $\forall x_a \in \CX_a$, $\foralmostall x_W \in \CX_W$: $g_b(x_a,x_{J\sm\{a\}},x_W) = g_b(x_a^0,x_{J\sm\{a\}},x_W)$;
    \item[(v)]  $\forall x_{J\sm\{a\}} \in \CX_{J\sm\{a\}}$, $\forall x_a \in \CX_a$: $P(g_b(x_a,x_{J\sm\{a\}},X_W)) = P(g_b(x_a^0,x_{J\sm\{a\}},X_W))$.
  \end{enumerate}
  Then (i) $\implies$ (ii) $\implies$ (iii) $\implies$ (iv) $\implies$ (v).
  If $M$ is $i\sigma$-faithful, then (v) $\implies$ (i).
\end{Prp}
\begin{proof}
  (i) $\implies$ (ii):
  Applying Lemma~\ref{lem:scm_structural_minimality_solution_function} with $B = \Anc^{G^+(M)}(b) \cap V$ and noting that
  $a \not\in \Anc^{G(M)}(b) \iff a \not\in \Anc^{G^+(M)}(b) \iff a \not\in \Pa^{G^+(M)}(B)$, we obtain that $g_b$ is constant in $x_a$.
  (ii) $\implies$ (iii) $\implies$ (iv) $\implies$ (v) is trivial.
  (v) $\implies$ (i) follows from the definition of $i\sigma$-faithfulness in combination with Remark~\ref{rem:scm_exogenous_input_causes_alt}.
\end{proof}
\Joris{
  UNIFICATION:
\begin{Prp}\label{prp:scm_causes}
  Let $M = \lt J, V, W, \CX, P, f \rt$ be a simple iSCM with causal graph $G(M)$.
  Let $a \in J \cup V$, $b \in V$. 
  Pick an arbitrary $x_{I_a}^0 \in \CX_{I_a}$ (for example, $x_{I_a}^0 = \star$ for $a \in V$).
  Write ``$\foralmostall x_w \in \CX_W$'' as a shorthand of ``for $P$-almost all $x_w \in \CX_W$''.
  Denote the solution function of $M_{\doit(I_a)}$ by $g$.
  Consider the properties:\footnote{Note that the ordering of the quantifiers matters! In particular, $\foralmostall x_W \forall x_V \implies \forall x_V \foralmostall x_W$, whereas the converse doesn't hold in general (but does hold if $\CX_V$ is discrete).}
  \begin{enumerate}[label=(\roman*)]
    \item[(i)]   $a \not\in \Anc^{G(M)}(b)$;
    \item[(ii)]  $\forall x_{J\sm\{I_a\}} \in \CX_{J\sm\{I_a\}}$, $\forall x_W \in \CX_W$, $\forall x_{I_a} \in \CX_{I_a}$: $g_b(x_{I_a},x_{J\sm\{I_a\}},x_W) = g_b(x_{I_a}^0,x_{J\sm\{I_a\}},x_W)$;
%    \item[(iii)] $\foralmostall x_W \in \CX_W$, $\forall x_{J\sm\{a\}} \in \CX_{J\sm\{a\}}$, $\forall x_a \in \CX_a$: $g_b(x_a,x_{J\sm\{a\}},x_W) = g_b(x_a^0,x_{J\sm\{a\}},x_W)$;
    \item[(iii)]  $\forall x_{J\sm\{I_a\}} \in \CX_{J\sm\{I_a\}}$, $\foralmostall x_W \in \CX_W$, $\forall x_{I_a} \in \CX_{I_a}$: $g_b(x_{I_a},x_{J\sm\{I_a\}},x_W) = g_b(x_{I_a}^0,x_{J\sm\{I_a\}},x_W)$;
    \item[(iv)]   $\forall x_{J\sm\{I_a\}} \in \CX_{J\sm\{I_a\}}$, $\forall x_{I_a} \in \CX_{I_a}$, $\foralmostall x_W \in \CX_W$: $g_b(x_{I_a},x_{J\sm\{I_a\}},x_W) = g_b(x_{I_a}^0,x_{J\sm\{I_a\}},x_W)$;
    \item[(v)]  $\forall x_{J\sm\{I_a\}} \in \CX_{J\sm\{I_a\}}$, $\forall x_{I_a} \in \CX_{I_a}$: $P(g_b(x_{I_a},x_{J\sm\{I_a\}},X_W)) = P(g_b(x_{I_a}^0,x_{J\sm\{I_a\}},X_W))$.
  \end{enumerate}
  Then (i) $\implies$ (ii) $\implies$ (iii) $\implies$ (iv) $\implies$ (v).
  If $M$ is $i\sigma$-faithful, then (v) $\implies$ (i).
\end{Prp}
\begin{proof}
  (i) $\implies$ (ii):
  Applying Lemma~\ref{lem:scm_structural_minimality_solution_function} to $M_{\doit(I_a)}$ with $B = \Anc^{G^+(M_{\doit(I_a)})}(b) \cap V$ and noting that
  $a \not\in \Anc^{G(M)}(b) \iff a \not\in \Anc^{G^+(M)}(b) \iff a \not\in \Anc^{G^+(M_{\doit(I_a)})}(b) \iff I_a \not\in \Pa^{G^+(M_{\doit(I_a)})}(B)$, we obtain that $g_b$ is constant in $x_{I_a}$.
  (ii) $\implies$ (iii) $\implies$ (iv) $\implies$ (v) is trivial.
  (v) $\implies$ (i) follows from the definition of $i\sigma$-faithfulness in combination with Remark~\ref{rem:scm_exogenous_input_causes_alt}.
\end{proof}
}
\begin{Lem}\label{lem:scm_structural_minimality_solution_function}
  Let $M = \lt J, V, W, \CX, P, f \rt$ be a simple iSCM with graph $G^+ := G^+(M)$.
  Let $B \subseteq V$.
  Denote the (unique) solution function of $M$ by $g$. %, and the (unique) solution function of the submodel $M^{[B]}$ by $g^{[B]}$.
  Then $g_B : \CX_J \times \CX_W \to \CX_B$ is constant in $(J \cup W) \sm \Pa^{G^+}(B)$:\footnote{We slightly abuse notation by
  identifying two functions with different domains.}
    $$\forall x \in \CX: \qquad g_B(x_J,x_W) = g_B(x_{(J \cup W) \cap \Pa^{G^+}(B)}).$$
%  and $g^{[B]} : \CX_J \times \CX_W \times \CX_{V \sm B} \to \CX_B$ is constant in $(J \cup W \cup (V \sm B)) \sm \Pa^{G^+}(B)$:
%    $$\forall x \in \CX: \qquad g^{[B]}(x_J,x_W,x_{V \sm B}) = g^{[B]}(x_{\Pa^{G^+}(B) \sm B}).$$
\end{Lem}
\begin{proof}
  One conceptually elegant way to prove this would be to apply the variation independence version of the global Markov property.
  Since we have not formulated this here, we will give a direct proof.

  First, note that for every $v \in V$, there exists a ``structurally minimal'' function 
  $\tilde{f}_v : \CX_{\Pa^{G^+}(v)} \to \CX_v$, a ``version'' of the causal mechanism for $v$ that only depends on the parents of $v$:
  $$x_v = f_v(x) \iff x_v = \tilde{f}_v(x_{\Pa^{G^+}(v)})$$
  for all $x \in \CX$. 
  This follows by induction from Definition~\ref{def:scm_parent}.
  Hence, for all $x \in \CX$,
  $$x_B = f_B(x) \iff x_B = \tilde{f}_B(x_{\Pa^{G^+}(B)}) = \tilde{f}_B(x_{\Pa^{G^+}(B) \sm B},x_{B \cap \Pa^{G^+}(B)}).$$
  Also, for all $x \in \CX$,
  $$x_B = f_B(x) \iff x_B = g^{[B]}(x_J,x_W,x_{V \sm B})$$
  where $g^{[B]}$ is the (unique) solution function of $M^{[B]}$.
  By the uniqueness of the solution function of $M^{[B]}$, then, the function $g^{[B]}$ must be constant in $x_k$ for $k \in (J \cup W \cup (V \sm B)) \sm (\Pa^{G^+}(B) \sm B)$.
  Furthermore, since for all $x \in \CX$, $x_B = g_B(x_J,x_W)$ implies $x_B = g^{[B]}(x_J,x_W,x_{V \sm B})$, also $g_B$ must be constant in $x_k$ for $k \in (J \cup W) \sm (\Pa^{G^+}(B) \sm B) = (J \cup W) \sm \Pa^{G^+}(B)$.
\end{proof}
\begin{Rem}\label{rem:scm_exogenous_input_causes_alt}
  Alternative equivalent formulations for some of the properties in Proposition~\ref{prp:scm_exogenous_input_causes} are:
  \begin{enumerate}
    \item[(i')] $a$ is not a cause of $b$ according to $G(M)$;
    \item[(i'')] $$b \isPerp_{G(M)} a \mid J \setminus \{a\};$$
  %    (follows from Proposition~\ref{prp:graph_cause_of})
    \item[(iv')] $a$ is not a rung-3 cause of $b$ according to $M$;
    \item[(v')] $a$ is not a rung-2 cause of $b$ according to $M$;
    \item[(v'')] $$X_b \Indep_{P_{M}} X_a \given X_{J\sm\{a\}};$$
  \end{enumerate}
\end{Rem}

\begin{Eg}\label{eg:scm_cause_of_iv_notimplies_iii}
  Consider the iSCM $M$ with exogenous input variable $X_a \in \RN$, endogenous variable $X_b \in \RN$, exogenous random variable $X_w \in \RN$, and structural equation
  $$X_b = \ind_{X_a}(X_w)$$
  and exogenous distribution $X_w \sim \C{N}(0,1)$.
  Then for all $x_a \in \RN$, $g(x_a,X_w) = 0$ a.s., so property (iv) in Proposition~\ref{prp:scm_exogenous_input_causes} holds.
  However, it's not the case that for almost all $X_w$, we have that $g_b(x_a,X_w) = g_b(x_a',X_w)$ for all $x_a,x_a'$, which means that property (iii) in Proposition~\ref{prp:scm_exogenous_input_causes} does not hold.
\end{Eg}

\begin{Eg}\label{eg:scm_cause_of_iii_notimplies_ii}
  Consider the iSCM $M$ with exogenous input variable $X_a \in \{0,1\}$, endogenous variable $X_b \in \{0,1\}$, exogenous random variable $X_w \in \{0,1\}$, and structural equations
  $$X_b = X_a \wedge X_w$$
  with $P(X_w=0) = 1$.
  Then $X_b = 0$ with probability 1, hence property (iii) in Proposition~\ref{prp:scm_exogenous_input_causes} holds.
  However, property (ii) in Proposition~\ref{prp:scm_exogenous_input_causes} does not hold.
\end{Eg}

\begin{Eg}\label{eg:scm_cause_of_ii_notimplies_i}
  Consider the iSCM $M$ with exogenous input variable $X_a \in \RN$, endogenous variables $X_c, X_d, X_b \in \RN$, exogenous random variable $X_w \in \RN$, and structural equations
  \begin{align*}
    X_c &= X_a, \\
    X_d &= -X_a \\
    X_b &= X_c + X_d + X_w 
  \end{align*}
  with $X_w \sim \C{N}(0,1)$.
  Then (i) in Proposition~\ref{prp:scm_exogenous_input_causes} does not hold, but (ii) in Proposition~\ref{prp:scm_exogenous_input_causes} holds.
\end{Eg}

\begin{Rem}\label{rem:scm_exogenous_input_causes_counterex}
  The reverse implications need not hold. Some counterexamples are: 
  \begin{itemize}
    \item (v) $\notimplies$ (iv): Example~\ref{eg:scm_cause_of_rung2_vs_rung3}.
    \item (iv) $\notimplies$ (iii): Example~\ref{eg:scm_cause_of_iv_notimplies_iii}.
    \item (iii) $\notimplies$ (ii): Example~\ref{eg:scm_cause_of_iii_notimplies_ii}.
    \item (ii) $\notimplies$ (i): Example~\ref{eg:scm_cause_of_ii_notimplies_i}.
  \end{itemize}
  Note that (iii) and (iv) are equivalent if $\CX_a$ is discrete.
  Note that (ii) and (iii) are equivalent if $\CX_W$ is discrete and the exogenous distribution is strictly positive.
\end{Rem}

\Joris{Some conjectures that I might prove some day:
\begin{itemize}
  \item Properties (i), (ii) and (iii) are not invariant under counterfactual equivalence.
  \item By redefining the graph of an iSCM in the spirit of \cite{BFPM21} (that is if $\forall X_J \foralmostall X_W \forall X_V$ a function is constant in some input then there is no edge),
    then (ii) is no longer relevant; we can then show that if all observed variables are discrete, the ancestral relations in the graph are identifiable from the counterfactual equivalence class.
\end{itemize}}

\JorisThoughts{
LEFTOVER REMARKS:
\begin{Eg}
  Take acyclic SCM $M$ defined by $\CX_a = \CX_b = \RN$, $\CX_w = \RN$, $X_w \sim \C{N}(0,1)$, structural equations $X_a = f_a(X_w) = X_w, X_b = f_b(X_a,X_w) = 0$.
  Take also acyclic SCM $\tilde{M}$ defined by $\CX_a = \CX_b = \RN$, $\CX_w = \RN$, $X_w \sim \C{N}(0,1)$, structural equations $X_a = \tilde{f}_a(X_w) = X_w, X_b = \tilde{f}_b(X_a,X_w) = \ind_{X_a}(X_w) X_a$.
  The solution function of $M$ is $g(X_w) = (X_w,0)$ while the solution function of $\tilde{M}$ is $\tilde{g}(X_w) = (X_w,X_w)$.
  This holds even though for all $x_a \in \CX_a$, $\tilde{f}_b(x_a,x_w) = 0$ for $P(X_w)$-almost every $x_w$.
  Note that under $\doit(X_a=x_a)$, $g^{\doit(x_a)}(X_w) = (x_a,0)$ and $\tilde{g}^{\doit(x_a)}(X_w) = (x_a,\ind_{x_a}(X_w) x_a) \stackrel{a.s.}{=} (x_a,0)$.
  
  This example shows that we cannot weaken the notion of equivalence in the AOS paper by swapping the ordering of the quantifiers, without losing the property that equivalent SCMs admit the same solutions. However, the proper ordering of the quantifiers with input nodes is still: $\forall X_J \foralmostall X_W \forall X_V$. 
  If we were to test whether $a$ causes $b$ by first turning $a$ into an input node, it might appear that we get the ordering $\forall X_a \foralmostall X_w \forall X_b$ in this example.
  However, turning $a$ into an input node means that its value can \emph{no longer} be determined by $X_w$, by the causal semantics of input nodes.
  Hence, it then indeed seems as if $X_a$ \emph{does not} cause $X_b$. Tricky!
  This suggests that if $X_a$ does not cause $X_b$ according to $M_{\doit(a)}$, it could still be that $X_a$ causes $X_b$ according to $M$.
  Well, this is perhaps not so surprising after all, and makes perfect sense. The notion $X_a$ causes $X_b$ according to $G_{\doit(a)}$ is weaker than the notion $X_a$ causes $X_b$ according to $G$.
\end{Eg}
}

For an endogenous variable $a$, we can apply the notions above to $M_{\doit(I_a)}$ with $I_a$ instead of $a$ (and then translate the statements back to $M$).
\begin{Cor}\label{cor:scm_endogenous_causes}
  Let $M = \lt J, V, W, \CX, P, f \rt$ be a simple iSCM with causal graph $G(M)$.
  Let $a,b \in V$. 
  Write ``$\foralmostall x_w \in \CX_W$'' as a shorthand of ``for $P$-almost all $x_w \in \CX_W$''.
  Denote the solution function of $M$ by $g$ and that of $M_{\doit(a)}$ by $g^{\doit(a)}$.
  Consider the properties:
  \begin{enumerate}[label=(\roman*)]
    \item[(i)]   $a \not\in \Anc^{G(M)}(b)$;
    \item[(ii)]  $\forall x_J \in \CX_J$, $\forall x_W \in \CX_W$, $\forall x_a \in \CX_a$:        $g_b^{\doit(a)}(x_a,x_W,x_J)    = g_b(x_W,x_J)$;
    \item[(iii)]  $\forall x_J \in \CX_J$, $\foralmostall x_W \in \CX_W$, $\forall x_a \in \CX_a$: $g_b^{\doit(a)}(x_a,x_W,x_J)    = g_b(x_W,x_J)$;
    \item[(iv)]   $\forall x_J \in \CX_J$, $\forall x_a \in \CX_a$, $\foralmostall x_W \in \CX_W$: $g_b^{\doit(a)}(x_a,x_W,x_J)    = g_b(x_W,x_J)$;
    \item[(v)]  $\forall x_J \in \CX_J$, $\forall x_a \in \CX_a$:                                $P(g_b^{\doit(a)}(x_a,X_W,x_J)) = P(g_b(X_W,x_J))$.
  \end{enumerate}
  Then (i) $\implies$ (ii) $\implies$ (iii) $\implies$ (iv) $\implies$ (v).
  If $M_{\doit(I_a)}$ is $i\sigma$-faithful, then (v) $\implies$ (i).
\end{Cor}
\begin{proof}
  Apply Proposition~\ref{prp:scm_exogenous_input_causes} to $M_{\doit(I_a)}$ with $I_a$ instead of $a$,
  and pick $a_0 = \star$, the value of $X_{I_a}$ that corresponds with ``no intervention on $a$''.
\end{proof}
\Joris{This is indeed now a special case of Proposition~\ref{prp:scm_causes}.}

\begin{Rem}\label{rem:scm_endogenous_causes_alt}
  Alternative formulations for some of the properties in Corollary~\ref{cor:scm_endogenous_causes} are:
  \begin{enumerate}
    \item[(i')] $a$ does not cause $b$ according to $G(M)$;
    \item[(i'')] $$b \isPerp_{G(M_{\doit(I_a)})} I_a \mid J;$$
    \item[(iv')] $a$ is not a rung-3 cause of $b$ according to $M$;
    \item[(v')] $a$ is not a rung-2 cause of $b$ according to $M$;
    \item[(v'')] $$X_b \Indep_{P_{M_{\doit(I_a)}}} X_{I_a} \given X_J.$$
  \end{enumerate}
\end{Rem}
\begin{Rem}\label{rem:scm_exogenous_input_causes_counterex}
  Similarly as in Remark~\ref{rem:scm_exogenous_input_causes_counterex}, the reverse implications need not hold.
  The counterexamples can be easily adapted to the case $a \in V$ by setting $X_a = x_a^0$ for some value $x_a^0 \in \CX_a$.
\end{Rem}

The reader may wonder why we chose to work with $M_{\doit(I_a)}$ rather than $M_{\doit(a)}$. 
The reason is that the observable Markov kernel may provide additional information beyond that provided by the interventional Markov kernels, as the following example shows.
\begin{Eg}\label{eg:scm_cause_of_include_observational_old}
  Consider the SCM $M$ with endogenous variables $X_a, X_b \in \{-1,+1\}$, exogenous random variable $X_w \in \{-1,1\}$, and structural equations:
  \begin{align*}
    X_a &= X_w, \\
    X_b &= X_a X_w,
  \end{align*}
  where $X_w \sim \C{U}(\{-1,1\})$.
  Note that $P_M(X_b = 1 \given \doit(X_a = x_a)) = \tfrac{1}{2}$ for $x_a \in \{-1,1\}$.
  However, $P_M(X_b = 1) = 1$.
  So we can not detect that $a$ causes $b$ from the pair of interventional Markov kernel $P_M(X_b \given \doit(X_a))$,
  but we can if we additionally compare with the observational distribution $P_M(X_b)$.
  Further, note that in contrast with Example~\ref{eg:scm_cause_of_rung2_vs_rung3} (which considers $M_{\doit(a)}$), here $a$ is a rung-2 cause of $b$ according to $M$.
  So this example shows that $a$ can be a rung-2 cause of $b$ according to $M$, while $a$ is not a rung-2 cause of $b$ according to $M_{\doit(a)}$ (which might be a little surprising at first sight).
  \Joris{Note: somehow this has become a duplicate of Example~\ref{ex:scm_cause_subtlety}?}
\end{Eg}

While we have given names to the notions (i), (iv) and (v), we consider notions (ii) and (iii) not as relevant (although they are useful as intermediate notions when deriving the logical relationships between the various notions).
Indeed, notion (ii) is too strong if we cannot intervene to set values of exogenous random variables, but are limited to drawing samples from $P(X_W)$.\footnote{This observation also suggests to change the definition of edges in the graph of an iSCM, analogously to how this is defined in \cite{BFPM21}. \Joris{This remark works better in the next section.}}
Notion (iii) is less relevant because it does not have an interpretation in terms of the data generating processes that we are considering:
it corresponds with first drawing a value for $X_W$ from $P(X_W)$, and afterwards deciding on a value for $X_a$. 
In practice, this would require the agent or experimenter to have access to the realization of $X_W$ when deciding on the value of $X_a$.
However, that is inconsistent with the modeling assumption that $X_W$ is latent for the agent.\footnote{One could also consider a setting in
which some, but not all, information about the value $x_W$ is accessible to the experimenter. For example, some exogenous random variables
could be considered ``observable'' and others ``latent'', while they would all be considered ``non-intervenable''. We will not consider this as 
it would complicate matters further.}

\Joris{Well, actually: if we assume all exgeonous random variables to be observable but non-intervenable, then notion (iii) does make sense
(assuming that the experimenter can observe and set to infinite precision).}

We have formalized the main principle of how we can learn about causal relations in the world: by actively changing some part of the world (choosing the intervention values independently) and observing the response of other parts of the world.
The independence assumption is key to distinguish mere correlation from causation.\footnote{As a less mathematical and more philosophical footnote: it is interesting to speculate about how this relates to the notion of a free will. 
If an agent is not convinced that it chose the intervention values independently of other past aspects of the world, it cannot validly perform this causal reasoning step. 
An agent without a free will to choose these values could therefore never conclude that its actions have a causal effect on the world, as it could also just be a puppet steered by higher powers, and any dependence it observes between its actions and aspects of the world could also be ascribed to confounding. 
So perhaps that is why evolution equipped us with the impression that we have a free will.}

\Joris{LEFTOVER REMARKS:

NOTE: notion (iii) seems irrelevant in the sense that it doesn't seem to have an interpretation in terms of a data generating process: 
we first sample $X_W$ and afterwards can decide on a value for $X_a$. 
Aha! If we intervene, we ARE NOT ALLOWED TO MAKE THE TARGET VALUE DEPEND ON $X_W$. Thus by intervening we loose such possible dependencies.
Moving certain forall-quantifiers to the right of the foralmostall-quantifier is a way to allow for such dependencies when intervening.
So we can look at two types of interventions: those where we first measure $X_W$ and then decide, and those where we cannot do this.
Interesting, the ``non-intervention'' is allowed to access $X_W$ when determining the outcome! So ``nature'' has larger power than the experimenter.
This implies that our notion of ``causes'' is allowed to be strictly weaker than that of equivalence of iSCMs, as the latter ALSO has to take into account non-interventions.
In other words, to explain the dependences that arise in ``natural'' experiments, we may need more edges in the graph than if we only want to explain dependences in experiments performed by the experimenter.
Since nature can pick values dependent on $X_W$, it can ``amplify'' the dependence on $X_a$ and also the $X_J$-conditional dependence on $X_a$.
Well, perhaps\dots even nature cannot pick $X_J$ dependent on $X_W$\dots
Aha, but it CAN pick values of the endogenous variables dependent on $X_W$, i.e., the values of $X_{V\sm \{a,b\}}$.
So we might look at the parent notion (with $a,b \in V$)
$a$ is a parent of $b$ iff $\forall X_J \foralmostall X_W \forall X_V: x_b = f_b(x_V,x_W,x_J) \iff x_b = f_b(x_{V\sm\{a\}},x_a^0,x_W,x_J)$.

(ii) also seems irrelevant if we cannot set the values of exogenous random nodes. With the `too strong' definition of edges in $G(M)$ we chose
here, (i) is natural, but if we weaken the definition of $G(M)$ also (i) will get weaker.
NOTE also that if we add mechanism changes, we can still have dependence on $X_W$! However, it seems that we still have to decide
on the mechanism before we sample the value of $X_W$, so this doesn't really change anything: it would still be incompatible with
the type of data generating processes we have in mind.}

\Joris{
For didactical purposes, we give here another proof that notion (i) implies (v) by using the do-calculus.

With the help of the do-calculus, we can connect the graphical notion (Definition~\ref{def:graph_cause_of}) and the rung-2 iSCM notion (Definition~\ref{def:scm_rung2_cause_of}) of causal relation.
The following proposition expresses formally that only if there is a directed path from $a$ to $b$ in $G(M)$, a hard intervention on $a$ can change the distribution of $b$ according to $M$. 
%With the help of the do-calculus, we can now tie this notion to practical procedures to deduce the presence of causal relations from the (intervened) Markov kernels of the iSCM.
%The following proposition expresses that only if $a$ causes $b$ according to $M$, a hard intervention on $a$ can change the distribution of $b$. 
%This further formalizes our intuitive notion of what it means for a variable to cause another variable.
%\begin{Prp}\label{prp:scm_causes_separation}
%  Let $M = \lt \Sin, \Send, \Srnd, \CX, P, f \rt$ be a simple iSCM with causal graph $G(M)$.
%  Let $a \in V \cup J$ and $b \in V$.
%  \begin{itemize}
%    \item For $a \in V$:
%      $$a \notin \Anc^{G(M)}(b) \implies P_{M}(X_b \given \doit(X_J), \doit(X_a)) = P_{M}(X_b \given \doit(X_J));$$
%    \item for $a \in J$:
%      $$a \notin \Anc^{G(M)}(b) \implies P_{M}(X_b \given \doit(X_J)) = P_{M}(X_b \given \doit(X_{J\setminus \{a\}}),\cancel{\doit(X_a)}).$$
%  \end{itemize}
%  Note: both cases can also be combined into a single statement:
%    $$a \notin \Anc^{G(M)}(b) \implies P_{M}(X_b \given \doit(X_{J \cup \{a\}}) = P_{M}(X_b \given \doit(X_{J\sm\{a\}})).$$
%\end{Prp}
\begin{Prp}\label{prp:scm_graph_cause_of}
  Let $M = \lt \Sin, \Send, \Srnd, \CX, P, f \rt$ be a simple iSCM with causal graph $G(M)$.
  Let $a \in V \cup J$ and $b \in V$.
  If $a$ does not cause $b$ according to $G(M)$, then $a$ is not a rung-2 cause of $b$ according to $M$.
\end{Prp}
\begin{proof}
  Denote $G = G(M)$. Assume that $a \notin \Anc^G(b)$.

  First, consider the case $a \in V$.
  By Proposition~\ref{prp:graph_cause_of},
  $$b \isPerp_{G_{\doit(I_a)}} I_a \mid J \setminus \{a\},$$
  where $I_a$ is an intervention node with target $a$ corresponding to an exogenous input variable modeling a hard intervention on $a$.
  We now apply rule 3 of the do-calculus (Theorem~\ref{thm:as-do-calculus-scms-simplified}) to obtain:
    $$P_{M}(X_b \given \doit(X_J), \doit(X_a)) = P_{M}(X_b \given \doit(X_J)).$$

  Second, consider the case $a \in J$. 
  By Proposition~\ref{prp:graph_cause_of},
    $$b \isPerp_{G_{\doit(a)}} a \mid J \setminus \{a\}.$$
  In that case we do not add an intervention node. 
  We now apply rule 4 of the do-calculus (Theorem~\ref{thm:as-do-calculus-scms-simplified}) to obtain:
    $$P_{M}(X_b \given \doit(X_J)) = P_{M}(X_b \given \doit(X_{J\setminus \{a\}}),\cancel{\doit(X_a)}).$$

%  \Joris{Patrick suggests to use the Markov property:
%  $$A \isPerp_{G(M)} B \given C \implies X_A \Indep_{P_{M}(X_V,X_W \given \doit(\XJ))} X_B \given X_C.$$
%  So in our case b)
%  $$X_b \Indep_{P_{M}(X_V,X_W \given \doit(\XJ))} X_a \given X_{J\sm \{a\}}.$$
%  Writing this out we get that there exists a Markov kernel $Q(X_b|X_{J\sm a})$ such that
%  $P_M(X_b,X_J | \doit(\XJ)) = Q(X_b|X_{J\sm a}) \otimes P_M(X_J|\doit(\XJ))$.
%  Hence
%  $P_M(X_b | \doit(\XJ)) = Q(X_b|X_{J\sm a})$.}

%  \Joris{Idea: could we shorten the proof by looking at the marginalized graph instead?}
\end{proof}
%This leads to a practical way of detecting the presence of a causal relation.
This gives sufficient conditions to identify the presence of (graphical) causal relations.
%\begin{Cor}\label{cor:scm_identify_causal_relation}
%  Let $M = \lt \Sin, \Send, \Srnd, \CX, P, f \rt$ be a simple iSCM.
%  Let $b \in V$.
%
%  \begin{itemize}
%    \item For $a \in V \cup W$ with $a \ne b$, if:
%  \begin{itemize}
%    \item there exist values $x_J \in \CX_J$, $x_a,x_a' \in \CX_a$ such that
%      $$P_{M}(X_b \given \doit(X_a=x_a),\doit(X_J=x_J)) \ne P_{M}(X_b \given \doit(X_a=x_a'),\doit(X_J=x_J)),$$
%    \item or there exist values $x_J \in \CX_J$, $x_a \in \CX_a$ such that
%  $$P_{M}(X_b \given \doit(X_a=x_a),\doit(X_J=x_J)) \ne P_{M}(X_b \given \doit(X_J=x_J)),$$
%  \end{itemize}
%  then $a$ is a cause of $b$ according to $M$.  
%    \item For $a \in J$, if there exist values $x_{J\setminus\{a\}} \in \CX_J$, $x_a,x_a' \in \CX_a$ 
%  such that
%      \begin{equation*}\begin{split}
%        & P_{M}(X_b \given \doit(X_a=x_a),\doit(X_{J\setminus \{a\}}=x_{J\setminus \{a\}})) \\
%        {} \ne {} & P_{M}(X_b \given \doit(X_a=x_a'),\doit(X_{J\setminus\{a\}}=x_{J\setminus\{a\}})),
%      \end{split}\end{equation*}
%  then $a$ is a cause of $b$ according to $M$.
%  \end{itemize}
%  Thus, in both cases, $a \in \Anc^{G(M)}(b)$.
%\end{Cor}
\begin{Cor}\label{cor:scm_identify_graphical_causal_relation}
  Let $M = \lt \Sin, \Send, \Srnd, \CX, P, f \rt$ be a simple iSCM.
  Let $b \in V$.

  For $a \in V \cup J$ with $a \ne b$, if there exist values $x_{J\sm\{a\}} \in \CX_{J\sm\{a\}}$, $x_a,x_a' \in \CX_a$ such that
    \begin{equation*}\begin{split}
      & P_{M}(X_b \given \doit(X_a=x_a),\doit(X_{J\setminus \{a\}}=x_{J\setminus \{a\}})) \\
      {} \ne {} & P_{M}(X_b \given \doit(X_a=x_a'),\doit(X_{J\setminus\{a\}}=x_{J\setminus\{a\}})),
    \end{split}\end{equation*}
  then $a$ is a rung-2 cause of $b$ according to $M$, and hence, $a$ is a cause of $b$ according to $G(M)$. %(that is, $a \in \Anc^{G(M)}(b)\sm\{b\}$).

  Also, for $a \in V$ with $a \ne b$, if there exist values $x_J \in \CX_J$, $x_a \in \CX_a$ such that
  $$P_{M}(X_b \given \doit(X_a=x_a),\doit(X_J=x_J)) \ne P_{M}(X_b \given \doit(X_J=x_J)),$$
  then $a$ is a rung-2 cause of $b$ according to $M$, and hence, $a$ is a cause of $b$ according to $G(M)$.  
\end{Cor}
We have seen that these conditions are sufficient, but not necessary.
}



\subsection{Causal Relations}

\subsubsection{Hierarchy of causal relations}

\Joris{Some conjectures that I might prove some day:
\begin{itemize}
  \item Properties (i), (ii) and (iii) are not invariant under counterfactual equivalence.
  \item By redefining the graph of an iSCM in the spirit of \cite{BFPM21} (that is if $\forall X_J \foralmostall X_W \forall X_V$ a function is constant in some input then there is no edge),
    then (ii) is no longer relevant; we can then show that if all observed variables are discrete, the ancestral relations in the graph are identifiable from the counterfactual equivalence class.
\end{itemize}}

\JorisThoughts{LEFTOVER REMARKS:
\begin{Eg}
  Take acyclic SCM $M$ defined by $\CX_a = \CX_b = \RN$, $\CX_w = \RN$, $X_w \sim \C{N}(0,1)$, structural equations $X_a = f_a(X_w) = X_w, X_b = f_b(X_a,X_w) = 0$.
  Take also acyclic SCM $\tilde{M}$ defined by $\CX_a = \CX_b = \RN$, $\CX_w = \RN$, $X_w \sim \C{N}(0,1)$, structural equations $X_a = \tilde{f}_a(X_w) = X_w, X_b = \tilde{f}_b(X_a,X_w) = \ind_{X_a}(X_w) X_a$.
  The solution function of $M$ is $g(X_w) = (X_w,0)$ while the solution function of $\tilde{M}$ is $\tilde{g}(X_w) = (X_w,X_w)$.
  This holds even though for all $x_a \in \CX_a$, $\tilde{f}_b(x_a,x_w) = 0$ for $P(X_w)$-almost every $x_w$.
  Note that under $\doit(X_a=x_a)$, $g^{\doit(x_a)}(X_w) = (x_a,0)$ and $\tilde{g}^{\doit(x_a)}(X_w) = (x_a,\ind_{x_a}(X_w) x_a) \stackrel{a.s.}{=} (x_a,0)$.
  
  This example shows that we cannot weaken the notion of equivalence in the AOS paper by swapping the ordering of the quantifiers, without losing the property that equivalent SCMs admit the same solutions. However, the proper ordering of the quantifiers with input nodes is still: $\forall X_J \foralmostall X_W \forall X_V$. 
  If we were to test whether $a$ causes $b$ by first turning $a$ into an input node, it might appear that we get the ordering $\forall X_a \foralmostall X_w \forall X_b$ in this example.
  However, turning $a$ into an input node means that its value can \emph{no longer} be determined by $X_w$, by the causal semantics of input nodes.
  Hence, it then indeed seems as if $X_a$ \emph{does not} cause $X_b$. Tricky!
  This suggests that if $X_a$ does not cause $X_b$ according to $M_{\doit(a)}$, it could still be that $X_a$ causes $X_b$ according to $M$.
  Well, this is perhaps not so surprising after all, and makes perfect sense. The notion $X_a$ causes $X_b$ according to $G_{\doit(a)}$ is weaker than the notion $X_a$ causes $X_b$ according to $G$.
\end{Eg}
}

\JorisThoughts{LEFTOVER REMARKS:

NOTE: notion (iii) seems irrelevant in the sense that it doesn't seem to have an interpretation in terms of a data generating process: 
we first sample $X_W$ and afterwards can decide on a value for $X_a$. 
Aha! If we intervene, we ARE NOT ALLOWED TO MAKE THE TARGET VALUE DEPEND ON $X_W$. Thus by intervening we loose such possible dependencies.
Moving certain forall-quantifiers to the right of the foralmostall-quantifier is a way to allow for such dependencies when intervening.
So we can look at two types of interventions: those where we first measure $X_W$ and then decide, and those where we cannot do this.
Interesting, the ``non-intervention'' is allowed to access $X_W$ when determining the outcome! So ``nature'' has larger power than the experimenter.
This implies that our notion of ``causes'' is allowed to be strictly weaker than that of equivalence of iSCMs, as the latter ALSO has to take into account non-interventions.
In other words, to explain the dependences that arise in ``natural'' experiments, we may need more edges in the graph than if we only want to explain dependences in experiments performed by the experimenter.
Since nature can pick values dependent on $X_W$, it can ``amplify'' the dependence on $X_a$ and also the $X_J$-conditional dependence on $X_a$.
Well, perhaps\dots even nature cannot pick $X_J$ dependent on $X_W$\dots
Aha, but it CAN pick values of the endogenous variables dependent on $X_W$, i.e., the values of $X_{V\sm \{a,b\}}$.
So we might look at the parent notion (with $a,b \in V$)
$a$ is a parent of $b$ iff $\forall X_J \foralmostall X_W \forall X_V: x_b = f_b(x_V,x_W,x_J) \iff x_b = f_b(x_{V\sm\{a\}},x_a^0,x_W,x_J)$.

(ii) also seems irrelevant if we cannot set the values of exogenous random nodes. With the `too strong' definition of edges in $G(M)$ we chose
here, (i) is natural, but if we weaken the definition of $G(M)$ also (i) will get weaker.
NOTE also that if we add mechanism changes, we can still have dependence on $X_W$! However, it seems that we still have to decide
on the mechanism before we sample the value of $X_W$, so this doesn't really change anything: it would still be incompatible with
the type of data generating processes we have in mind.}

For didactical purposes, we give here another proof that notion (i) implies (v) by using the do-calculus.

With the help of the do-calculus, we can connect the graphical notion (Definition~\ref{def:graph_cause_of}) and the rung-2 iSCM notion (Definition~\ref{def:scm_rung2_cause_of}) of causal relation.
The following proposition expresses formally that only if there is a directed path from $a$ to $b$ in $G(M)$, a hard intervention on $a$ can change the distribution of $b$ according to $M$. 
%With the help of the do-calculus, we can now tie this notion to practical procedures to deduce the presence of causal relations from the (intervened) Markov kernels of the iSCM.
%The following proposition expresses that only if $a$ causes $b$ according to $M$, a hard intervention on $a$ can change the distribution of $b$. 
%This further formalizes our intuitive notion of what it means for a variable to cause another variable.
%\begin{Prp}\label{prp:scm_causes_separation}
%  Let $M = \lt \Sin, \Send, \Srnd, \CX, P, f \rt$ be a simple iSCM with causal graph $G(M)$.
%  Let $a \in V \cup J$ and $b \in V$.
%  \begin{itemize}
%    \item For $a \in V$:
%      $$a \notin \Anc^{G(M)}(b) \implies P_{M}(X_b \given \doit(X_J), \doit(X_a)) = P_{M}(X_b \given \doit(X_J));$$
%    \item for $a \in J$:
%      $$a \notin \Anc^{G(M)}(b) \implies P_{M}(X_b \given \doit(X_J)) = P_{M}(X_b \given \doit(X_{J\setminus \{a\}}),\cancel{\doit(X_a)}).$$
%  \end{itemize}
%  Note: both cases can also be combined into a single statement:
%    $$a \notin \Anc^{G(M)}(b) \implies P_{M}(X_b \given \doit(X_{J \cup \{a\}}) = P_{M}(X_b \given \doit(X_{J\sm\{a\}})).$$
%\end{Prp}
\begin{Prp}\label{prp:scm_graph_cause_of}
  Let $M = \lt \Sin, \Send, \Srnd, \CX, P, f \rt$ be a simple iSCM with causal graph $G(M)$.
  Let $a \in V \cup J$ and $b \in V$.
  If $a$ does not cause $b$ according to $G(M)$, then $a$ is not a rung-2 cause of $b$ according to $M$.
\end{Prp}
\begin{proof}
  Denote $G = G(M)$. Assume that $a \notin \Anc^G(b)$.

  First, consider the case $a \in V$.
  By Proposition~\ref{prp:graph_cause_of},
  $$b \isPerp_{G_{\doit(I_a)}} I_a \mid J \setminus \{a\},$$
  where $I_a$ is an intervention node with target $a$ corresponding to an exogenous input variable modeling a hard intervention on $a$.
  We now apply rule 3 of the do-calculus (Theorem~\ref{thm:as-do-calculus-scms-simplified}) to obtain:
    $$P_{M}(X_b \given \doit(X_J), \doit(X_a)) = P_{M}(X_b \given \doit(X_J)).$$

  Second, consider the case $a \in J$. 
  By Proposition~\ref{prp:graph_cause_of},
    $$b \isPerp_{G_{\doit(a)}} a \mid J \setminus \{a\}.$$
  In that case we do not add an intervention node. 
  We now apply rule 4 of the do-calculus (Theorem~\ref{thm:as-do-calculus-scms-simplified}) to obtain:
    $$P_{M}(X_b \given \doit(X_J)) = P_{M}(X_b \given \doit(X_{J\setminus \{a\}}),\cancel{\doit(X_a)}).$$

%  \Joris{Patrick suggests to use the Markov property:
%  $$A \isPerp_{G(M)} B \given C \implies X_A \Indep_{P_{M}(X_V,X_W \given \doit(\XJ))} X_B \given X_C.$$
%  So in our case b)
%  $$X_b \Indep_{P_{M}(X_V,X_W \given \doit(\XJ))} X_a \given X_{J\sm \{a\}}.$$
%  Writing this out we get that there exists a Markov kernel $Q(X_b|X_{J\sm a})$ such that
%  $P_M(X_b,X_J | \doit(\XJ)) = Q(X_b|X_{J\sm a}) \otimes P_M(X_J|\doit(\XJ))$.
%  Hence
%  $P_M(X_b | \doit(\XJ)) = Q(X_b|X_{J\sm a})$.

%  \Joris{Idea: could we shorten the proof by looking at the marginalized graph instead?}
\end{proof}
%This leads to a practical way of detecting the presence of a causal relation.
This gives sufficient conditions to identify the presence of (graphical) causal relations.
%\begin{Cor}\label{cor:scm_identify_causal_relation}
%  Let $M = \lt \Sin, \Send, \Srnd, \CX, P, f \rt$ be a simple iSCM.
%  Let $b \in V$.
%
%  \begin{itemize}
%    \item For $a \in V \cup W$ with $a \ne b$, if:
%  \begin{itemize}
%    \item there exist values $x_J \in \CX_J$, $x_a,x_a' \in \CX_a$ such that
%      $$P_{M}(X_b \given \doit(X_a=x_a),\doit(X_J=x_J)) \ne P_{M}(X_b \given \doit(X_a=x_a'),\doit(X_J=x_J)),$$
%    \item or there exist values $x_J \in \CX_J$, $x_a \in \CX_a$ such that
%  $$P_{M}(X_b \given \doit(X_a=x_a),\doit(X_J=x_J)) \ne P_{M}(X_b \given \doit(X_J=x_J)),$$
%  \end{itemize}
%  then $a$ is a cause of $b$ according to $M$.  
%    \item For $a \in J$, if there exist values $x_{J\setminus\{a\}} \in \CX_J$, $x_a,x_a' \in \CX_a$ 
%  such that
%      \begin{equation*}\begin{split}
%        & P_{M}(X_b \given \doit(X_a=x_a),\doit(X_{J\setminus \{a\}}=x_{J\setminus \{a\}})) \\
%        {} \ne {} & P_{M}(X_b \given \doit(X_a=x_a'),\doit(X_{J\setminus\{a\}}=x_{J\setminus\{a\}})),
%      \end{split}\end{equation*}
%  then $a$ is a cause of $b$ according to $M$.
%  \end{itemize}
%  Thus, in both cases, $a \in \Anc^{G(M)}(b)$.
%\end{Cor}
\begin{Cor}\label{cor:scm_identify_graphical_causal_relation}
  Let $M = \lt \Sin, \Send, \Srnd, \CX, P, f \rt$ be a simple iSCM.
  Let $b \in V$.

  For $a \in V \cup J$ with $a \ne b$, if there exist values $x_{J\sm\{a\}} \in \CX_{J\sm\{a\}}$, $x_a,x_a' \in \CX_a$ such that
    \begin{equation*}\begin{split}
      & P_{M}(X_b \given \doit(X_a=x_a),\doit(X_{J\setminus \{a\}}=x_{J\setminus \{a\}})) \\
      {} \ne {} & P_{M}(X_b \given \doit(X_a=x_a'),\doit(X_{J\setminus\{a\}}=x_{J\setminus\{a\}})),
    \end{split}\end{equation*}
  then $a$ is a rung-2 cause of $b$ according to $M$, and hence, $a$ is a cause of $b$ according to $G(M)$. %(that is, $a \in \Anc^{G(M)}(b)\sm\{b\}$).

  Also, for $a \in V$ with $a \ne b$, if there exist values $x_J \in \CX_J$, $x_a \in \CX_a$ such that
  $$P_{M}(X_b \given \doit(X_a=x_a),\doit(X_J=x_J)) \ne P_{M}(X_b \given \doit(X_J=x_J)),$$
  then $a$ is a rung-2 cause of $b$ according to $M$, and hence, $a$ is a cause of $b$ according to $G(M)$.  
\end{Cor}
We have seen that these conditions are sufficient, but not necessary.

\subsection{Direct Causal Relations}

%\begin{Def}\label{def:direct_cause_of}
%  Let $M = \lt \Sin, \Send, \Srnd, \CX, P, f \rt$ be a simple iSCM with causal graph $G = G(M)$.
%  Let $a \in J \cup V$ and $b \in V$.
%  We say that \emph{$a$ is a direct cause of $b$ (w.r.t.\ $V \cup J$) according to $M$} if $a \in \Pa^{G}(b) \sm \{b\}$.
%\end{Def}
%\begin{Prp}\label{prp:direct_cause_of}
%  Let $M = \lt \Sin, \Send, \Srnd, \CX, P, f \rt$ be a simple iSCM.
%  For $a \in J \cup V$, $b \in V$, $a$ is a direct cause of $b$ w.r.t.\ $J \cup V$ according to $M$ if and only if $a$ is a cause of $b$ according to $M^{[b]}$.
%\end{Prp}
Applying Corollary~\ref{cor:scm_identify_causal_relation} to the submodel $M^{[\{a,b\}\cap V]}$ gives similar conditions to identify the presence of direct causal relations.
%\begin{Cor}\label{cor:scm_identify_direct_causal_relation}
%  Let $M = \lt \Sin, \Send, \Srnd, \CX, P, f \rt$ be a simple iSCM.
%  Let $b \in V$.
%
%  \begin{itemize}
%    \item For $a \in V \cup W$ with $a \ne b$, if:
%  \begin{itemize}
%    \item there exist values $x_J \in \CX_J$, $x_{V\setminus\{a,b\}} \in \CX_{V\setminus\{a,b\}}$, $x_a,x_a' \in \CX_a$ such that
%      \begin{equation*}\begin{split}
%        & P_{M}(X_b \given \doit(X_a=x_a),\doit(X_{V\setminus\{a,b\}}=x_{V\setminus\{a,b\}}),\doit(X_J=x_J)) \\
%        {} \ne {} & P_{M}(X_b \given \doit(X_a=x_a'),\doit(X_{V\setminus\{a,b\}}=x_{V\setminus\{a,b\}}),\doit(X_J=x_J)),
%      \end{split}\end{equation*}
%    \item or there exist values $x_J \in \CX_J$, $x_{V\setminus\{a,b\}} \in \CX_{V\setminus\{a,b\}}$, $x_a \in \CX_a$ such that
%      \begin{equation*}\begin{split}
%        & P_{M}(X_b \given \doit(X_a=x_a),\doit(X_{V\setminus\{a,b\}}=x_{V\setminus\{a,b\}}),\doit(X_J=x_J)) \\
%        {} \ne {} & P_{M}(X_b \given \doit(X_{V\setminus\{a,b\}}=x_{V\setminus\{a,b\}}),\doit(X_J=x_J)),
%      \end{split}\end{equation*}
%  \end{itemize}
%  then $a$ is a direct cause of $b$ w.r.t.\ $V \cup J \cup W$ according to $M$.
%    \item For $a \in J$, if there exist values $x_{J\setminus\{a\}} \in \CX_J$, $x_{V\setminus\{a,b\}} \in \CX_{V\setminus\{a,b\}}$, $x_a,x_a' \in \CX_a$ 
%  such that
%      \begin{equation*}\begin{split}
%        & P_{M}(X_b \given \doit(X_a=x_a),\doit(X_{V\setminus\{a,b\}}=x_{V\setminus\{a,b\}}),\doit(X_{J\setminus \{a\}}=x_{J\setminus \{a\}})) \\
%        {} \ne {} & P_{M}(X_b \given \doit(X_a=x_a'),\doit(X_{V\setminus\{a,b\}}=x_{V\setminus\{a,b\}}),\doit(X_{J\setminus\{a\}}=x_{J\setminus\{a\}})),
%      \end{split}\end{equation*}
%  then $a$ is a direct cause of $b$ w.r.t\ $V \cup J \cup W$ according to $M$.
%  \end{itemize}
%  Thus, in both cases, $a \in \Pa^{G(M)}(b)$.
%\end{Cor}
\begin{Cor}\label{cor:scm_identify_direct_causal_relation}
  Let $M = \lt \Sin, \Send, \Srnd, \CX, P, f \rt$ be a simple iSCM.
  Let $b \in V$.

  For $a \in V \cup J$ with $a \ne b$, if there exist values $x_{J\setminus\{a\}} \in \CX_{J\sm\{a\}}$, $x_{V\setminus\{a,b\}} \in \CX_{V\setminus\{a,b\}}$, $x_a,x_a' \in \CX_a$ such that
      \begin{equation*}\begin{split}
        & P_{M}(X_b \given \doit(X_a=x_a),\doit(X_{V\setminus\{a,b\}}=x_{V\setminus\{a,b\}}),\doit(X_{J\setminus \{a\}}=x_{J\setminus \{a\}})) \\
        {} \ne {} & P_{M}(X_b \given \doit(X_a=x_a'),\doit(X_{V\setminus\{a,b\}}=x_{V\setminus\{a,b\}}),\doit(X_{J\setminus\{a\}}=x_{J\setminus\{a\}})),
      \end{split}\end{equation*}
  then $a$ is a direct cause of $b$ w.r.t.\ $V \cup J$ according to $M$, and hence, according to $G(M)$.

  Also, for $a \in V$ with $a \ne b$, if there exist values $x_J \in \CX_J$, $x_{V\setminus\{a,b\}} \in \CX_{V\setminus\{a,b\}}$, $x_a \in \CX_a$ such that
      \begin{equation*}\begin{split}
        & P_{M}(X_b \given \doit(X_a=x_a),\doit(X_{V\setminus\{a,b\}}=x_{V\setminus\{a,b\}}),\doit(X_J=x_J)) \\
        {} \ne {} & P_{M}(X_b \given \doit(X_{V\setminus\{a,b\}}=x_{V\setminus\{a,b\}}),\doit(X_J=x_J)),
      \end{split}\end{equation*}
  then $a$ is a direct cause of $b$ w.r.t.\ $V \cup J$ according to $M$, and hence, according to $G(M)$.
%  Thus in both cases, $a \in \Pa^{G(M)}(b)$.
\end{Cor}
%\begin{proof}
%  Apply Corollary~\ref{cor:scm_identify_causal_relation} to the intervened iSCM 
%  $$M^{[b]} = \lt J \cup V \setminus \{b\}, \{b\}, \Srnd, \CX, P, f_{\{b\}} \rt,$$
%  and note that $a \in \Anc^{G(M^{[b]})}(b) \implies a \in \Pa^{G(M)}(b) \cup \{b\}$.
%\end{proof}

\JorisThoughts{NOTE: notion (iii) don't seem right, as it allows us to pick $x_a,x_a'$ depending on $X_W$. 
So if we define $j \tuh v$ not in $G(M)$ iff $g_v(x_J,x_W) = g_v(x_j^0,x_{J\sm\{j\}},x_W)$ for all $x_{J\sm\{j\}}$, for almost all $x_W$.
Well, that is for a single node. For multiple nodes, we should define, for $a \in J$:
\begin{center}
  $a \tuh b$ not in $G(M)$ for $a\in J, b \in V$ iff $x_b = f_b(x_J,x_V,x_W) \iff x_b = f_b(x_{J\sm a},x_a^0,x_V,x_W)$ 
$\forall x_{J\sm\{a\}}, \foralmostall x_W, \forall x_V$.
\end{center}
Applying this to $M_{\doit(V\sm\{a,b\}),\doit(I_a)}$ gives:
\begin{center}
  $a \tuh b$ in $G(M)$ for $a\in V, b \in V$ iff $x_b = f_b(x_J,x_V,x_W) \iff x_b = f_b(x_J,x_a^0,x_b,x_{V\sm\{a,b\}},x_W)$ 
$\forall x_J, \forall x_{V\sm\{a,b\}}, \foralmostall x_W, \forall x_a, \forall x_b$.
In words: if the solutions of the SE for $x_b$ do not depend on the value of $x_a$ when $\doit(I_a), \doit(V\sm\{a,b\})$.
\end{center}
This fits a general pattern of the quantifiers: for all exogenous inputs, for almost all exogenous randoms, for all endogenous ones.
PERHAPS: should we also $\doit(I_{V\sm\{a,b\}})$ to have more options? No! Then we don't block indirect effects.

Can we now prove that equivalent iSCMs have the same graph?
}
\JorisThoughts{
\begin{Con}
  If we had defined the directed edges in the graph of an iSCM as follows, they would have been counterfactually identifiable,
  that is, (i) would hold iff (v) would hold, and there would be no need to consider properties (ii), (iii) and (iv) anymore.

  For getting direct effects, we consider $M_{\doit(V\sm\{a,b\})}$.
  First, consider the edges $a \tuh b$ with $a \in J, b \in V$:
  \begin{center}
    for $a \in J, b\in V$: $a \tuh b$ in $G(M)$ if (by definition):
    $$\forall x_{J\sm\{a\}}, \forall x_{V\sm\{b\}}, \foralmostall x_W, \forall x_a, \forall x_b : x_b = f_b(x_J,x_V,x_W) \iff x_b = f_b(x_{J\sm a},x_a^0,x_V,x_W).$$
  \end{center}
  OR PERHAPS:
  \begin{center}
    for $a \in J, b\in V$: $a \tuh b$ in $G(M)$ if (by definition):
    $$\forall x_{J\sm\{a\}}, \forall x_{V\sm\{b\}}, \forall x_a, \foralmostall x_W, \forall x_b : x_b = f_b(x_J,x_V,x_W) \iff x_b = f_b(x_{J\sm a},x_a^0,x_V,x_W).$$
  \end{center}
  To define the edges $a \tuh b$ with $a, b \in V$, we apply this after $\doit(I_a)$:
  \begin{center}
    for $a,b \in V$: $a \tuh b$ in $G(M)$ if (by definition):
    $$\forall x_J, \forall x_{V\sm\{a,b\}}, \foralmostall x_W, \forall x_a, \forall x_b: x_b = f_b(x_J,x_V,x_W) \iff x_b = f_b(x_J,x_a^0,x_b,x_{V\sm\{a,b\}},x_W).$$
  \end{center}
  OR PERHAPS:
  \begin{center}
    for $a,b \in V$: $a \tuh b$ in $G(M)$ if (by definition):
    $$\forall x_J, \forall x_{V\sm\{a,b\}}, \forall x_a, \foralmostall x_W, \forall x_b : x_b = f_b(x_J,x_V,x_W) \iff x_b = f_b(x_J,x_a^0,x_V,x_W).$$
  \end{center}
  OR PERHAPS we should substitute $x_a^0$ by $g_a^{[a,b]}(x_{V\sm\{a,b\}},x_J,x_W)$?

  This captures the notion that the solutions of the structural equation for $x_b$ do not depend on the value of $x_a$ when $\doit(I_a), \doit(V\sm\{a,b\})$.

  Hm, the conjecture seems false. It appears that in order to get counterfactual identifiability of the graph, we need to be able to conclude from
    $$\forall x_J, \forall x_{V\sm\{a,b\}}, \forall x_a, \foralmostall x_W, \forall x_b: x_b = f_b(x_J,x_V,x_W) \iff x_b = f_b(x_J,x_a^0,x_b,x_{V\sm\{a,b\}},x_W)$$
  that $a \tuh b \notin G(M)$, i.e.~ with the proposed definition:
    $$\forall x_J, \forall x_{V\sm\{a,b\}}, \foralmostall x_W, \forall x_a, \forall x_b: x_b = f_b(x_J,x_V,x_W) \iff x_b = f_b(x_J,x_a^0,x_b,x_{V\sm\{a,b\}},x_W).$$
  This works if $x_a$ is discrete. Perhaps we could also pose this as an additional requirement for continuous iSCMs?
  The latter may be an option indeed. Without it, we can just have the type of counterexample we already identified in the notafiability note.
  So perhaps: only allow causal mechanisms for which we can pull each of the endogenous variables to the left of the foralmostall quantifier?
\end{Con}
}
\JorisThoughts{
  Nope. Salvaging Stephan's proof would require weakening the parent-ship relation to:
  $a$ is NOT parent of $b$ according to $M$ iff there exists a $\tilde{f}_b$ such that
  for all $x_{\sm b}$, for $P$-almost all $e$, for all $x_b$:
  $$x_b = f_b(x,e) \iff x_b = \tilde{f}_b(x_{\sm a},e)$$
  Or perhaps: iff there exists a $\tilde{f}_b$ such that
  for all $x_{\sm \{a,b\}}$, for $P$-almost all $e$, for all $x_{a,b}$:
  $$x_b = f_b(x,e) \iff x_b = \tilde{f}_b(x_{\sm a},e)$$
  Apparently, the first formulation came from analyzing the removed result from the AOS paper.
  The second formulation came from sticking closer to the ``$\foralmostall x_W \forall x_V$''
  point of departure, and from recognizing that by non-intervening on $x_a$ we end up picking
  a value of $x_a$ that depends on $x_W$, and hence we get more freedom in picking the
  null set of $x_W$ if we have $\forall x_a$ inner to $\foralmostall x_W$.
  This might just work!

  However, I believe that in the acyclic setting and for all endogenous variables discrete, 
  we can also salvage Stephan's proof.
  An interesting counterexample to consider is
  $$f_b(x_a,e) = \delta_{x_a,e}$$
  with $E \sim \C{N}(0,1)$. Even though for each $x_a$, $f_b(x_a,E) = 0$ $P(E)$-a.s.,
  there is no $P(E)$-null set such that $f_b(x_a,E) = 0$ for all $x_a$ and all $e$ in the
  null set complement.
  Could we consider this $f_b$ equivalent to the null function?
  Let's look at cycles. Consider
  $$\begin{aligned}
    x_a &= f_a(x_b,e_a) = \delta_{x_b,e_a} + \alpha x_b + e_a \\
    x_b &= f_b(x_a,e_b) = \delta_{x_a,e_b} + \beta x_a + e_b
  \end{aligned}$$
  Is this equivalent to 
  $$\begin{aligned}
    x_a &= \tilde{f}_a(x_b,e_a) = \alpha x_b + e_a \\
    x_b &= \tilde{f}_b(x_a,e_b) = \beta x_a + e_b
  \end{aligned}$$
  Could the ``jumps'' (which would be almost surely absent in the acyclic case) be amplified until a non-measure zero discontinuity?
%  a = \alpha b + e
%  b = \beta a + f
%  a = \alpha (\beta a + f) + e = \alpha \beta a + \alpha f + e
%  (1 - \alpha \beta) a = \alpha f + e
%  a = (\alpha f + e) / (1 - \alpha \beta)
%  b = (\beta e + f) / (1 - \alpha \beta)

%  a = \alpha b + e + (b==e)
%  b = \beta a + f + (a==f)
%  a = \alpha (\beta a + f + (a==f)) + e + (b==e)
%  (1 - \alpha \beta) a = \alpha f + e + \alpha (a==f) + (b==e)

%  f = (\alpha f + e) / (1 - \alpha \beta)
%  (1 - \alpha \beta) f = \alpha f + e
%  (1 - \alpha \beta - \alpha) f = e
%  f = e / (1 - \alpha \beta - \alpha)

 % So as long as e!=b and f!=e/(1-\alpha\beta-\alpha) the solution coincides with that of the non-jumpy structural equations.

  So if we can prove that the new parentship notion allows us to first remove all non-parents, and then solving a set
  of structural equations, up to a null set, we may be in business! Or, in other words, if the equivalence relation
  corresponding to the new parentship relation preserves solutions, we are in business.

  Can we do this in the acyclic setting?
  Equivalence notion of AOS paper: for almost all $e$, for all $x$:
  $$x_i = f_i(x,e) \iff x_i = \tilde{f}_i(x,e)$$
  That is, in terms of solution functions of the submodels $[i]$, $g^{M^{[i]}}$ and $g^{\tilde{M}^{[i]}}$,
  if for almost all $e$,
  $$\forall x_{\sm i}: g^{M^{[i]}}(x_{\sm i},e) = g^{\tilde{M}^{[i]}}(x_{\sm i},e)$$

  The idea is to change this into: for all $x_{\sm i}$, up to a $P(E)$-nullset $N_{x_{\sm i}}$,
  $$g^{M^{[i]}}(x_{\sm i},e) = g^{\tilde{M}^{[i]}}(x_{\sm i},e).$$
  That is, we turned around the ordering of the quantifiers.
  We can also write this as: for all $x_{\sm i}$, for almost all $e$:
  $$\forall x_i: x_i = f_i(x_i,x_{\sm i},e) \iff x_i = \tilde{f}_i(x_i,x_{\sm i},e)$$

  Can we show that in that case they admit the same solutions?
  Let $(X,E)$ be a solution of $M$. Then $X=f(X,E)$ a.s..
  That is, $P(\omega : X(\omega)=f(X(\omega),E(\omega)))=1$.
  }


\subsection{Common Causes}

\begin{Prp}\label{prp:scm_graph_common_cause}
  Let $M = \lt \Sin, \Send, \Srnd, \CX, P, f \rt$ be a simple iSCM with causal graph $G(M)$.
  Let $a, b \in V$ with $a \ne b$, and $c \in V \cup J$.
  If $c$ is a common cause of $a$ and $b$ according to $M$, then $c$ is a common cause of $a$ and $b$ according to $G(M)$.
\end{Prp}
\begin{proof}
  Assume that $c$ is a common cause of $a$ and $b$ according to $M$.
  By definition, then, $c$ is a cause of $a$ according to $M_{\doit(b)}$ and $c$ is a cause of $b$ according to $M_{\doit(a)}$.
  Since $c$ causes $a$ according to $M_{\doit(b)}$, by Proposition~\ref{prp:scm_graph_cause_of} applied to $M_{\doit(b)}$, $c$ causes $a$ according to $G(M_{\doit(b)}) = G(M)_{\doit(b)}$.
  Similarly, $c$ causes $b$ according to $G(M)_{\doit(a)}$.
  Hence, $c \in \Anc^{G(M)_{\doit(b)}}(a)$ and $c \in \Anc^{G(M)_{\doit(a)}}(b)$.
  By Proposition~\ref{prp:bifurcations_alternative}, there exists a bifutcation with source $c$ in $G$ between $a$ and $b$.
  Hence $c$ is a common cause of $a$ and $b$ according to $G(M)$.
\end{proof}
 
To test whether some variable is a common cause of two other variables, we can simply apply  Corollary~\ref{cor:scm_identify_causal_relation} for detecting causal relations (after appropriate interventions), making use of Proposition~\ref{prp:bifurcations_alternative} to reexpress the existence of a bifurcation with source in terms of certain ancestral relations after performing certain interventions. 
%We will not write this out in detail here.

\begin{Cor}\label{cor:scm_identify_common_cause}
Let $M = \lt \Sin, \Send, \Srnd, \CX, P, f \rt$ be a simple iSCM with causal graph $G(M)$.
Let $a, b \in V$ with $a \ne b$ and $c \in V \cup J \sm \{a,b\}$.
%If $c \in \Anc^{G_{\doit(b)}}(a) \sm \{a\}$ and $c \in \Anc^{G_{\doit(a)}}(b) \sm \{b\}$ then $c$ is a common cause of $a$ and $b$ according to $M$.

If there exist values $x_a \in \CX_a$, $x_{J\sm\{c\}} \in \CX_{J\sm\{c\}}$, $x_c,x_c' \in \CX_c$ such that
  \begin{equation*}\begin{split}
    & P_M(X_b \given \doit(X_c=x_c),\doit(X_{J\setminus \{c\}}=x_{J\setminus \{c\}}),\doit(X_a=x_a)) \\
    {} \ne {} & P_M(X_b \given \doit(X_c=x_c'),\doit(X_{J\setminus\{c\}}=x_{J\setminus\{c\}}),\doit(X_a=x_a)),
  \end{split}\end{equation*}
and if there exist (possibly different) values $x_b \in \CX_b$, $x_{J\sm\{c\}} \in \CX_{J\sm\{c\}}$, $x_c,x_c' \in \CX_c$ such that
  \begin{equation*}\begin{split}
    & P_M(X_a \given \doit(X_c=x_c),\doit(X_{J\setminus \{c\}}=x_{J\setminus \{c\}}),\doit(X_b=x_b)) \\
    {} \ne {} & P_M(X_a \given \doit(X_c=x_c'),\doit(X_{J\setminus\{c\}}=x_{J\setminus\{c\}}),\doit(X_b=x_b)),
  \end{split}\end{equation*}
then $c$ is a common cause of $a$ and $b$ according to $M$, and hence, according to $G(M)$.
\end{Cor}
\begin{proof}
Combine Proposition~\ref{prp:bifurcations_alternative} with Corollary~\ref{cor:scm_identify_causal_relation}.
\end{proof}  
We used here only the first sufficient condition in Corollary~\ref{cor:scm_identify_causal_relation}. One can similarly make use of the second sufficient condition in Corollary~\ref{cor:scm_identify_causal_relation}, possibly in combination with the first, to derive additional sufficient conditions for detecting common causes.
Note that these sufficient conditions are also not necessary.

%\Joris{Another convention could be: all output nodes without incoming edges are assumed to be exogenous random nodes.
%One can assume this WLOG. It seems to make sense for a confounder. I.e., an endogenous variable without parents
%should not be considered a confounder.}

\subsection{Confounding}

\subsubsection{Confounding (graphical notion)}

\Joris{The following proposition may be redundant!}
\begin{Prp}\label{prp:graph_unconfounded1}
  Let $G$ be a CDMG with input nodes $J$ and output nodes $V$.
  Let $a,b \in V$ be distinct output nodes.
  If $b$ does not cause $a$ according to $G$, and $a$ and $b$ are unconfounded according to $G$, then
  $$b \isPerp_{G_{\doit(I_a)}} I_a \given \{a\} \cup J.$$
\end{Prp}
\begin{proof}
Assume that $b \notin \Anc^G(a)$ and every bifurcation in $G$ between $a$ and $b$ has a source in $J$.

We proceed by contradiction.
Suppose that there exists a walk in $G_{\doit(I_a)}$ between $b$ and $I_a \cup J$ in $G_{\doit(I_a)}$ that is $(J \cup \{a\})$-$\sigma$-open and such that all colliders lie in $J \cup \{a\}$.
It cannot contain nodes from $J$, as such a node would either be an end node of the walk ($\sigma$-blocking it)
or a non-collider node pointing only to nodes in another strongly connected component of $G_{\doit(I_a)}$
($\sigma$-blocking the walk).
Hence it must be of the form $I_a \tuh a \cdots b$. 
If it is not of the form $I_a \tuh a \hus \cdots b$ with only a single occurrence of $a$, we can modify it to turn it into this form, while keeping it $(J \cup \{a\})$-$\sigma$-open.
Indeed, suppose it is of the form $I_a \tuh \cdots a \tuh c \cdots b$ (with $a$ the right-most $a$, that is, the occurrence of $a$ closest to the end point $b$, and possibly $c=b$) then $c$ must be in the same strongly connected component as $a$, and hence we can take $I_a \tuh a \hut \cdots \hut c \cdots b$ by replacing the subwalk between the left-most $a$ and $c$ by a directed path $a \hut \dots \hut c$ through nodes in the strongly connected component of $a$.
If it is of the form $I_a \tuh \cdots a \hus c \cdots b$ (with $a$ the right-most $a$, and possibly $c=b$), we can take $I_a \tuh a \hus c \cdots b$ by skipping the subwalk between the left-most $a$ and the right-most $a$.
If this path contains $b$ more than once, we can shorten it further so that it contains $b$ exactly once.

So there exists a walk of the form $I_a \tuh a \hus \cdots b$ in $G_{\doit(I_a)}$ that is $(J \cup \{a\})$-$\sigma$-open, contains $a$ and $b$ only once, contains no other colliders than $a$, and the subwalk between $a$ and $b$ is not a directed walk from $b$ to $a$.
Hence it must be of the form $I_a \tuh a \hut \cdots \hut c \tuh \cdots \tuh b$ (with $c \in V$) or
$I_a \tuh a \huh c \tuh \cdots \tuh b$ or $I_a \tuh a \hut \dots \hut c_1 \huh c_2 \tuh \dots \tuh b$.
As it contains $a$ and $b$ only once, we conclude the existence of a bifurcation in $G_{\doit(I_a)}$ between $a$ and $b$, either with source $c \in V$ or without source.
%Therefore, we can replace the subpath between $a$ and the right-most collider $c$ on the path by a directed path from that collider $c$ to $a$. 
%Note that this directed path cannot contain $b$, and cannot pass through any of the nodes in $\C{O}\setminus\{a,b\}$. 
%We thereby obtain a $\sigma$-open walk of the form $I_a \to a \oto b$ or $I_a \to a \ot \cdots \ot c \ot \cdots b$ or $I_a \to a \ot \cdots \ot c \oto \cdots b$. 
%We thereby obtain a $(J \cup \{a\})$-$\sigma$-open walk of the form $I_a \to a \oto \cdots b$ or $I_a \to a \ot \cdots b$, 
%where the part between $a$ and $b$ contains no colliders and is not a directed walk from $b$ to $a$. Hence it must be a bifurcation.
Contradiction.
Hence $b \isPerp_{G_{\doit(I_a)}} I_a \given \{a\} \cup J$.

\Joris{Alternative proof using graphical marginalization: consider $\tilde{G} := G^{\sm(V \sm \{a,b\})}$.
Consider $\tilde{G}_{\doit(I_a)}$. Since $b \notin \Anc^G(a)$, also $b \notin \Anc^{\tilde{G}}(a)$.
Because $a,b$ are the only output nodes in $\tilde{G}$, this means that $b \tuh a$ is not in $\tilde{G}$.
As there exists no bifurcation between $a$ and $b$ in $G$, there exists no bifurcation bewteen $a$ and $b$ in $\tilde{G}$.
Because $a,b$ are the only output nodes in $\tilde{G}$, this means that $b \huh a$ is not in $\tilde{G}$, and there is no $j \in J$ such that $b \hut j \tuh a$ in $\tilde{G}$.
So the only possible edge between $a$ and $b$ in $\tilde{G}$ is $a \tuh b$. 
%The only other possible edges in $\tilde{G}$ are of the form $j \tuh a$ and $j \tuh b$ for $j \in J$ (but not both $b \hut j \tuh a$).
Consider a path between $b$ and $I_a \cup J$ in $\tilde{G}_{\doit(I_a)}$. 
Suppose it is $\{a\} \cup J$-$\sigma$-open.
It cannot contain a node in $J$.
Hence it must be of the form $I_a \tuh a \sus b$.
As the only edge between $a$ and $b$ is $a \tuh b$, it can only be the path $I_a \tuh a \tuh b$.
But that is a contradiction, as that path is not $\{a\} \cup J$-$\sigma$-open.
Hence $b \isPerp_{\tilde{G}_{\doit(I_a)}} I_a \given \{a\} \cup J$.
Because $\tilde{G}_{\doit(I_a)} = (G_{\doit(I_a)})^{\sm (V \sm \{a,b\})}$, also $b \isPerp_{G_{\doit(I_a)}} I_a \given \{a\} \cup J$.
}
\end{proof}

\subsubsection{Confounding (counterfactual notion)}
  
\subsubsection{Confounding (potential outcomes notion)}

\JorisTODO{Exchangeability
might appear at first sight to be a weaker assumption than $a,b$ having no common causes according to some simple iSCM,
since it only requires a joint distribution on the various potential outcomes to be defined, and does not
presuppose the existence of an underlying iSCM that induces these distributions via the twin operation. 
However, perhaps the two notions can still be regarded as equivalent provided that $|\CX_a|$ is finite or countably infinite, 
as one can just construct a corresponding simple iSCM (with exogenous random variable $X_a$, one other independent exogenous
random variable, and endogenous variable $X_b$ such that $P_M(X_b \given \doit(X_a=x_a)) \sim X_{b'}^{\doit(X_{a'}=x_a)}$.
Here we need that $P(X_{b'}^{\doit(X_{a'}=x_a)})$ is a Markov kernel $\CX_a \dshto \CX_b$, in particular, that this is
measurably indexed by $x_a$.
Another option is to consider a simple iSCM with exogenous random variable $X_a$, exogenous input variable $X_{a'}$, and
endogenous variable $X_{b'}$.
See also \cite{RobinsRichardson2010}.}
  
%\Joris{Another convention could be: all output nodes without incoming edges are assumed to be exogenous random nodes.
%One can assume this WLOG. It seems to make sense for a confounder. I.e., an endogenous variable without parents
%should not be considered a confounder.}

%\Joris{We should perhaps define a canonical form of simple iSCMs. This form assumes:
%\begin{enumerate}
%  \item All factors in the exogenous random distribution are non-degenerate;
%  \item All endogenous nodes without parents have been replaced by exogenous random nodes.
%\end{enumerate}


\subsubsection{Confounding (interventional notion)}

\begin{Rem}\label{rem:cyclic_no_confounding_bias}
How to define ``confounding bias'' in cyclic simple iSCMs is an open research problem.
\Joris{In the acyclic case, if $i$ and $j$ are not confounded, we have the identity
      $$p(X_i,X_j \given \doit(X_{\C{O}\setminus\{i,j\}})) = 
       p(X_i \given \doit(X_j),\doit(X_{\C{O}\setminus\{i,j\}})) 
       p(X_j \given \doit(X_i),\doit(X_{\C{O}\setminus\{i,j\}}))$$

      In the unconfounded real-valued cyclic case, I believe I derived the equality
      $$\left(1 - \frac{\partial f_X}{\partial y}\frac{\partial f_Y}{\partial x}\right) p(X=x \given \doit(Y=y)) p(Y=y \given \doit(X=x)) = p(x,y)$$
      with $f_X(y,e_X) = F^{-1}_{X \given \doit(Y=y)}$ and $f_Y(x,e_Y) = F^{-1}_{Y \given \doit(X=x)}$ the inverse cdf's
      of the interventional distributions. In the acyclic case this gives the identity above.
      Note that $p(Y = y \given \doit(X=x))$ completely fixes $f_Y$ in $Y = f_Y(X,E_Y)$ up to the marginal
      distribution of $E_Y$ (which we can choose WLOG), $p(X = x \given \doit(Y=y))$ completely fixes $f_X$
      in $X = f_X(Y,E_X)$ up to the marginal distribution of $E_X$. In the unconfounded case, this fixes the joint
      distribution of $(E_X,E_Y)$ (by independence), and hence it fixes the joint density $p(X,Y)$. That is, even
      in the unconfounded cyclic case, the two interventional distributions completely determine the joint density.
      If there is confounding, this no longer holds (generically). This argument goes through for smooth enough 
      densities and functions and real-valued variables.

      From wikipedia (see also appendix on measure theory):
      $$\int_\Omega h \circ g d\mu = \int_{g(\Omega)} h dg_* \mu$$
      or, for $g : \Omega \to \RN^n$ a $C^1$ diffeomorphism with $\Omega \subseteq \RN^n$ an open subset:
      $$\int_\Omega h d\mu = \int_{g(\Omega)} h \circ g^{-1} dg_* \mu = \int_{g(\Omega)} h \circ g^{-1} |\det D_x T^{-1}| \,dm(x)$$
      with the Radon-Nikodym derivative w.r.t.\ Lebesgue measure given by:
      $$\frac{dg_*\lambda}{d\lambda}(x) = |\det D_x T^{-1}|$$

      Idea: if we represent $P(X \given \doit(Y))$ by a deterministic function, and also $P(Y \given \doit(X))$, then we can take these functions as the structural equations for an unconfounded iSCM, and read off the implied observational distribution via the push-forward of its solution function. The only question is: does the solution function exist, is it unique, and is it measurable?
}
\end{Rem}

\subsection{Direct confounding}

\begin{Def}\label{def:graph_directly_confounded}
  Let $M = \lt \Sin, \Send, \Srnd, \CX, P, f \rt$ be a simple iSCM with causal graph $G(M)$.
  Let $a,b \in V$ be distinct endogenous variables.
  If $a \huh b$ is present in $G(M)$ then we say that \emph{$a$ and $b$ are directly confounded (w.r.t.\ $V \cup J$) according to $M$}.
\end{Def}
\begin{Rem}\label{rem:graph_directly_confounded}
  In terms of the graph $G^+(M)$ of $M$, this holds if and only if there exists $w \in W$ such that $a \hut w \tuh b$ in $G^+(M)$.
\end{Rem}
\begin{Def}\label{def:scm_directly_confounded}
Let $M = \lt \Sin, \Send, \Srnd, \CX, P, f \rt$ be a simple iSCM.
Let $a, b \in V$ with $a \ne b$.
If for all values $x_J \in \CX_J, x_V \in \CX_V$, $f_a(x_J,X_W,x_V) \Indep f_b(x_J,X_W,x_V)$ then we say that \emph{$a$ and $b$ are ``directly'' unconfounded according to $M$}.
  \Joris{This might work in the acyclic case, but in the cyclic case we really have to solve $M^{[a,b]}$.}
\end{Def}
\begin{Prp}\label{prp:directly_confounded}
  Let $M = \lt \Sin, \Send, \Srnd, \CX, P, f \rt$ be a simple iSCM with causal graph $G(M)$.
  For $a,b \in V$, $a$ and $b$ are directly confounded (w.r.t.\ $V \cup J$) according to $M$ if and only if $a$ and $b$ are confounded according to $M^{[\{a,b\}]}$.
  \Joris{??? update!}
\end{Prp}

\begin{Prp}\label{prp:scm_unconfoundedness}
Let $M = \lt \Sin, \Send, \Srnd, \CX, P, f \rt$ be a simple iSCM.
Let $a \ne b \in V$.
If $b \notin \Anc^{G(M)}(a)$ and $a$ and $b$ are not confounded according to $M$, then
  $$P_{M}(X_b \given \doit(X_a), \doit(X_{V\setminus\{a,b\}}), \doit(X_J)) = P_{M}(X_b \given X_a, \doit(X_{V\setminus\{a,b\}}), \doit(X_J)) \qquad \mu_a\text{-a.s.}$$
for any reference measure $\mu_a$ on $\CX_a$ with $\mu_a \ll P_M(X_a \given \doit(X_J)) \ll \mu_a$.
\end{Prp}
%The contra-positive can be used as a sufficient condition for detecting that two nodes are confounded according to $M$.

For the simpler representation, we can intervene on all variables in $V\sm\{a,b\}$ and then only a few possible edges remain: $a \tuh b$, $b \tuh a$, $a \huh b$. 
If $a \huh b$ isn't there, we get the separation assumption and can draw the conclusion of no ``direct'' confounding bias (if we intervene to fix all other variables).

This leads to a criterion to detect a bidirected edge!

\subsection{Modeling implications}

Obviously (?), if $a$ and $b$ have a common cause then they are confounded.
But that is not the only possibility.
Two endogenous variables can also be confounded in case of latent selection bias, that is, if only a specific subset of the samples (statistical units, individuals, \dots) ends up in the data set because of some filtering process. 
\Joris{TODO: give examples}
\begin{comment}
Importantly: without common causes (and without a reverse causal relation), there can be no bias due to confounding.\footnote{Another bias is still possible: selection bias. This may happen if a subset of the samples (statistical units, individuals, \dots) do not end up in the data set because of some filtering process. 
In other words, even if
  $$P_{M}(X_b \given \doit(X_a), \doit(X_J)) = P_{M}(X_b \given X_a, \doit(X_J)) \qquad \mu_a\text{-a.s.}$$
we need not have
  $$P_{M}(X_b \given S \in \xi_S, \doit(X_a), \doit(X_J)) = P_{M}(X_b \given S \in \xi_S, X_a, \doit(X_J)) \qquad \mu_a\text{-a.s.}$$
for some $S \subseteq V \cup W$ and measurable subset $\xi_S \subseteq \CX_S$.
We plan to add more discussion on this to a future version of these lecture notes.}
\end{comment}


%Note that exogenous input nodes in $J$ are never called confounders, by definition.\footnote{The reason is that they do not lead to ``spurious dependences'' as long as one does not mix distributions for different values of $X_J$. 
%\Joris{How to deal with constant exogenous random nodes, or constant endogenous nodes without parents? We might want to not consider those as confounders. Also, the Markov property could be extended: by substituting the constant values in their children nodes (perhaps recursively) we can get rid of these nodes, apply the Markov property to the result, and we can freely add the removed nodes in any of the spots of any separation statement.}}

\subsubsection{Some remaining thoughts}
  \Joris{
    Some thoughts:
    \begin{itemize}
      \item How should we define ``confounded'' without referring to an iSCM or graph, but instead a set of Markov kernels corresponding to different data generating processes?
      \item There are still two interesting remaining notions for ``confounded according to an iSCM'' to consider that are invariant under interventional equivalence: 
        (i) there exists an interventionally equivalent iSCM for which $a,b$ are rung-3 unconfounded according to Definition~\ref{def:scm_rung3_unconfounded}; (ii) if we give each endogenous variable private independent copies of the exogenous random variables, it remains interventionally equivalent with the original iSCM. 
        It could be that the two notions are equivalent. Clearly, (ii) implies (i). Does (i) imply (ii)? 
      \item Example 0.3.14 in marginal\_scm\_causal\_inference.pdf sharpens Example~\ref{eg:interventionally_equivalent_but_different_graphs} by turning it into a faithful example with the same curious property (some directed edges not identifiable from interventional equivalence class).
      \item Example 0.3.18 in marginal\_scm\_causal\_inference.pdf shows that Corollary~\ref{cor:scm_identify_causal_relation} is not a necessary condition for identifiability of causal relations.
      \item Example 0.3.11 in marginal\_scm\_causal\_inference.pdf is worth adding (it's about stronger sufficient conditions to detect causal relations: sometimes we can detect the joint effect of $a$ on $b_1,b_2$ even though we cannot detect the effect of $a$ on $b_1$ separately, nor of $a$ on $b_2$).
      \item Proposition 0.3.3 in marginal\_scm\_causal\_inference.pdf might be worth writing down formally (it is about the equivalence of several notions of ``is a direct cause of'' for DAGs).
      \item Interestingly, Definition~\ref{def:graph_unconfounded} and Remark~\ref{rem:graph_unconfounded} can remain unaltered even when introducing latent selection bias.
        The confusing part is that the CDMG by itself does not express which variables are conditioned on.
        So we'd need to extend the notion of CDMG to allow for conditioning (what did the 'C' stand for, again?).
        We would like to represent such data generating processes graphically and more quantitatively.
        One can then see that there are more ways to end up with confounding than just marginalization.
        If any path is open, conditional on the filtering steps, between $a$ and $b$ that is not a directed path, then we call it confounded.
        No, that also doesn't work: in $a \tuh s \hut b$, are $a$ and $b$ confounded given $s=1$?
        Wait: in order to be able to distinguish ``confounding'' from ``pure causation'' between two nodes $a,b$, the nodes need to have a causal interpretation!
        Hence, if latent selection keeps the causal interpretation of $a,b$ intact, and we do not consider any directed path between $a$ and $b$, then the question is whether there are remaining open paths (after conditioning on $S$) between $a$ and $b$.
        So that would be a promising graphical criterion for (CDMG,selection nodes) representations of data generating processes.
    \item For testing no confounding bias according to $M$, could we make an independent copy of $X_W$ for each $v \in V$, and check whether the modified $M$ would be interventionally equivalent w.r.t.\ $\{a,b\}$ to the original $M$?
    \item We could attempt to introduce an ``absolute'' notion of unconfoundedness by saying that $a$ and $b$ are unconfounded if there exists an iSCM $M$ that reproduces $P(X_a \given \doit(X_b))$, $P(X_b \given \doit(X_a))$ and $P(X_a,X_b)$, such that according to $M$, $a$ and $b$ are unconfounded.
      This is related to our old definition 0.3.22 in marginal\_scm\_causal\_inference.pdf
      We also had a result (proposition 0.3.23) that connects that definition with the notion of ``no confounding bias''.
      Note that this ``absolute'' notion would still refer to a class of causal models.
      It would be nice if we can show that this property can be expressed purely in terms of properties of $P(X_a \given \doit(X_b))$, $P(X_b \given \doit(X_a))$ and $P(X_a,X_b)$.
      \item It actually seems that the SWIG approach to defining unconfoundedness also may work in case of cycles: one can write $Y(x) \Indep X$ even if $X \tuh Y$ and $X \hut Y$.
      One does have to cycle-split then, though. 
      The question is whether this is always realistic: the confounder that determines $X$ a split second before we intervene $\doit(x)$ and then wait for equilibration before we measure $Y(x)$ should not change over time. 
      But perhaps: if we already assumed that the same iSCM can model data generating processes leading to $X$ and $Y(x)$ both accurately, then why not also the above data generating process?
      Aha: the assumption is stronger! For the first we just assume that the distribution of the confounder doesn't change, for the second that its value doesn't even change!
      To test this, we could measure $X$ at different points in time and see whether its value is constant. 
      We could also do this under all possible interventions on its parents.
      This would give us quite some confidence in the latent values staying the same as well (but no guarantee!).
      Aha: it is the consistency assumption: if we assume that $Y(X) = Y$ then it means that if we measure $X$, then a little later intervene to set $\doit(X=X)$, then wait until equilibration, and we get the same value $Y(X) = Y$, then it is compatible with what the simple iSCM models. However, if in the real world, the value differs, e.g., because some background factor changes over time and $Y(X(t),U(t+1)) \ne Y(U(t))$, then the system violates the consistency assumption (and hence cannot be modelled well with a simple iSCM). By the way, this holds both for simple (cyclic) iSCMs and for acyclic iSCMs, but for cyclic iSCMs the assumption seems more restrictive because after intervening we need to give the system some time to equilibrate, during which the value of the confounder may change. 
      Note that even if $U(t)$ and $U(t+1)$ differ, they can be dependent, and then we can still use $Y(x) \Indep X$ as a test for confounding in the cyclic case!
      However, if at every $t$ an independent value $U(t)$ is drawn, we cannot detect the ``instantaneous'' confounding anymore.
      
      Aha. Perhaps this is what is meant by ``individal'' causality? An individual corresponds with a value of the exogenous random variables, and this is assumed invariant over time, so we can do many observations/interventions on the same object?

      This all, however, sounds more like ``how can we detect confounding'' than ``how should we define confounding''?
      Perhaps we can define it indirectly: if there exists no iSCM with exogenous inputs $J$, endogenous variables $a,b$ and without a bidirected edge $a \huh b$ that reproduces $P(X_a,X_b\given\doit(X_J))$, $P(X_a\given\doit(X_J),\doit(X_a))$ and $P(X_b\given\doit(X_J),\doit(X_b))$.
    \item Example 0.4.5 in marginal\_scm\_causal\_inference.pdf might be nice to add here.
  \end{itemize}
}
  \begin{comment}
  If there exists a $w \in W$ such that $w$ is a parent of both $a$ and $b$ according to $M_{\{a,b\}}$ then we say that \emph{$a$ and $b$ are confounded according to $M$}.
  \Joris{This seems a bit too strong: we might want to exclude degenerate (deterministic) $w$.}
  \Joris{Or perhaps: if we take $M_{\{a,b\}}$ and then take each $w \in W$ that is a common parent of $a$ and $b$ and split it into $w_a,w_b$ that we take independent, do we then get an iSCM that is interventionally equivalent to $M$ w.r.t $\{a,b\}$? Or perhaps even counterfactually?}
  \Joris{Indeed, if we just merge two independent separate causes of $a$ and $b$ into one variable, it looks like a confounder, but it actually doesn't introduce any dependencies between the two.}
  \Joris{This also relates to the ambiguity in the representation of the exogenous distribution. Perhaps we should at this point deal with it, once and for all. That is, make the notion of confounder independent of the chosen representation of the exogenous random variables? That would be cool! This would give us something like ``confounded according to an equivalence class of iSCMs''.}
  \Joris{Or, even more drastically: take $M_{\{a,b\}}$; if there exists any iSCM with endogenous variables $a,b$ and exogenous input variables $J$ (but arbitrary exogenous variables) that has no edge $a \huh b$ and is interventionally equivalent to $M_{\{a,b\}}$, then $a$ and $b$ are unconfounded according to $M$.}
  \Joris{This definition would also cover the cyclic (simple) iSCM case. It would also still allow the practical test 'causation is correlation' for the acyclic case.  Basically, this would amount to asking for the existence of an unconfounded iSCM in its interventional equivalence class.}
  \Joris{This definition seems adequate if the iSCM has only two endogenous variables ($a$ and $b$), but it's not clear whether it's appropriate otherwise.
  Indeed, we could imagine that (within the interventional equivalence class) other variables provide information to detect confounding.
  So we could also ask whether there is any other iSCM $\tilde{M}$ that is interventionally equivalent to $M$ and that has no common exogenous random parent of $a,b$ and not even an endogenous common cause of $a,b$; hm, does that even make sense?}
  \end{comment}

\subsection{Abstraction through Marginalization}

\Joris{Note also that the following notation is almost never used in this section! It seems more in line with the iSCM notation to write it as $G_O$ instead.
However, it is used quite a bit in the FCI chapter, and there it appears useful to remind us of the presence of the input nodes.
Note that the observable graph, currently written as $G_{O|J}$, \emph{also} marginalizes out nodes in $W$! This cannot be done in the iSCM currently,
so we could at best use $G(M_{O\cap V})^{\setminus (W\sm O)}$ instead. Note that $G(M_{O\cap V})^{\setminus (W\sm O)} \subseteq (G(M))^{\setminus((V \cup W) \setminus O)} =: G_{O|J}$.
I should investigate whether we only need the notation for $O \subseteq V$ (that is, all $W$ are by default considered unobservable).
By the way, a graph is not observable, so this is a bad name! Perhaps ``marginal graph'' instead?
Suppose we also allow for conditioning. Then ``marginal graph'' is a bad name. We could call it the ``effective graph'' perhaps in that case?}

We often deal with the situation that only part of the variables are observed or observable, and others are latent.
\begin{Not}
For a simple iSCM $M$, let $O \subseteq V \cup W$ be the set of observable variables,
where we will always assume $J$ to be observable as well.
We refer to the marginalized graph
$$G_{O|J}(M) := (G(M))^{\setminus((V \cup W) \setminus O)}$$
as the \emph{observable} graph of $M$. It has input nodes $J$ and output nodes $O$.
The corresponding observable Markov kernel is denoted 
  $$P_{M}(X_O \mid \doit(X_J)),$$
and the observable intervened Markov kernel resulting from a hard intervention targeting $T \subseteq O$ is 
  $$P_{M}(X_{O\setminus T} \mid \doit(X_{J\cup T})) := P_{M_{\doit(T)}}(X_{O\sm T} \mid \doit(X_{J\cup T})).$$
\end{Not}

Remember that Remark~\ref{rem:marg_preserves_ancestors_bifurcations_acyclicity} stated that graphical marginalization
preserves ancestral relations and bifurcations, while 
Lemma~\ref{lem:stability_separation_marginalization} stated that graphical marginalization preserves
$id$-separation statements and $i\sigma$-separation statements concerning observable variables.
These properties have powerful consequences:
\begin{Rem}
For many causal reasoning purposes, it suffices to know the observable graph $G_{O|J}(M)$ rather than $G(M)$.
In particular, the observable graph contains all information concerning:
  \begin{enumerate}
    \item which pairs of observable variables are causally related according to $M$ (c.f.\ Definition~\ref{def:cause_of});
    \item which pairs of observable variables have a common cause according to $M$ (c.f.\ Definition~\ref{def:common_cause_of});
    \item the $id$-separation and $i\sigma$-separation statements between observable variables according to $M$ (c.f.\ Definition~\ref{def:sigma_separation}).
  \end{enumerate}
\end{Rem}
In addition, the probabilistic marginalization preserves the notion of conditional independence. That is,
for $A, B, C \subseteq J \cup O$, we have
  $$X_A \Indep_{P_{M}(X_V,X_W \mid X_J)} X_B \mid X_C \iff X_A \Indep_{P_{M}(X_O \mid X_J)} X_B \mid X_C.$$
This means that for applying the global Markov property on observable variables (and hence the do-calculus 
and adjustment criteria), it suffices to know the observable graph $G_{O|J}(M)$ and the observable (intervened) 
Markov kernels rather than the full graph $G(M)$ and the full Markov kernels associated to $M$.
\textbf{In other words: these properties taken together are very powerful, as they allow us to abstract away irrelevant details when performing causal reasoning.}

\subsection{Abstraction through Conditioning}

\Joris{This is experimental still! One thing I run into now is that even though we can easily define the term ``confounded w.r.t.\ $V$'', it is not so easy as with the common cause to define ``confounded w.r.t.\ $O$'' for some observed subset $O \subseteq V$. The problem is that the iSCM doesn't represent the filtering mechanisms. So we should say ``confounded w.r.t.\ $O$ according to $M$ and filters $\{S_i \in \xi_i\}$, with $S_i \in V$, $\xi_i \subseteq \CX_i$. Here we could even distinguish singleton $\xi_i$ from non-singleton $\xi_i$. This becomes a bit awkward.

We could formally define a ``Selection iSCM'' as 
a tuple consisting of an iSCM and a finite list of conditioning events of the form $f_s(X_V,X_W,X_J) = 1$ with $f_s$ a binary function. 
One could think of introducing dummy endogenous variables with structural equations $S_s = f_s(X)$ but they are non-causal.

The conjecture is that any selection iSCM with latent non-causal exogenous random variables can be reduced to a non-selection iSCM with latent non-causal exogenous random variables.
This would explain that the class of iSCMs with latent non-causal exogenous random variables is a convenient class to work with, as the assumed properties can be made WLOG
(actually it increases the generality compared to the class of iSCMs with causal exogenous random variables).
}

A common convention in the iSCM literature is that exogenous random variables are ``non-causal'' in the sense that they are not allowed to be intervention targets.
This formal restriction (interventions can only target endogenous variables) can be exploited in modeling situations by allowing the modeler to be less specific about certain modeling details.
In other words: if one only cares about interventional equivalence w.r.t.\ the endogenous variables, there is still a lot of freedom in how one chooses the exogenous random variables and their distribution.
This freedom is increased further if exogenous random variables are considered latent (another often used convention in the iSCM literature).

If we do not give a causal interpretation to the exogenous random variables, it turns out that in addition to accounting for ``bias'' or ``spurious dependence'' due to a latent common cause, they can account for ``selection bias'' due to a selection step or filtering mechanism in the past. We will flesh out this interpretation in this subsection.

  \Joris{In the definition of confounding: How to deal with constant exogenous random nodes, or constant endogenous nodes without parents? We might want to not consider those as confounders. Also, the Markov property could be extended: by substituting the constant values in their children nodes (perhaps recursively) we can get rid of these nodes, apply the Markov property to the result, and we can freely add the removed nodes in any of the spots of any separation statement.
  Aha! If we consider exogenous random variables as non-causal, we can change the definition of $G(M)$ slightly by demanding that the parent-relationship of $w$ for $a,b$ holds $\Prb(X_w)$-a.s., instead of surely. Then we automatically obtain that the edges $a \hut w \tuh b$ disappear from $G(M)$ if $W$ is degenerate, because that just means $\Prb(X_w)$-a.s.\ constant, and that means it won't be a parent. Aha, this also all suggests that for a consistent treatment, we will have to change all ``for all $x_W$'' quantifiers into ``for almost all $x_W$'' quantifiers. Dashing a node then implies switching from forall to for-almost-all.}

\Joris{We could try showing that confounding may stem from any open path between $a$ and $b$ that's not a directed path? 
This however sounds as if we need to discuss the explicit representation of filtering steps.
Hm, it is in a way what the next proposition is showing, except that we need to add the explicit conditioning. 
}


\subsection{Reichenbach's Principle of Common Cause}

The following exercise illustrates this.
\begin{Exc}
  Reichenbach's \emph{Principle of Common Cause} states that if two events are dependent,
  then one must cause the other or the events must have a common cause (or any combination of these three possibilities).

  We can make this precise for simple iSCMs in the following way. Assume that $M$
  is a simple iSCM with two observed endogenous variables $X,Y$ (and possibly other
  latent variables as well, but no exogenous input nodes). Prove that 
  $X \nIndep_{P_{M}(X,Y)} Y$
  implies that $X \tuh Y$,
  $X \hut Y$ or $X \huh Y$ in $G_{\{X,Y\}}(M)$.

  Thus, either $X$ causes $Y$ according to $M$, 
  or $Y$ causes $X$ according to $M$,
  or $X$ and $Y$ have a common cause according to $M$.
\end{Exc}
